% Dedicated current SOURCE export; historical G26/G29 exports remain outside this cohort.
\section*{NeuroTranslation Definitions and Lookup Index}
\phantomsection\label{programme-definitions}
\addcontentsline{toc}{section}{NeuroTranslation Definitions and Lookup Index}
\noindent Canonical registry revision 34; 200 entries; SHA-256 \texttt{3fffa79b\allowbreak{}c9f404ca\allowbreak{}a7d4910e\allowbreak{}9e37de5a\allowbreak{}97371fdc\allowbreak{}1ad11bc5\allowbreak{}4155d309\allowbreak{}ec5f5394}.\par
\noindent A measurement, a computer state and a claim about biology need different words. These entries explain the current programme meanings. A proposed name does not establish a working interface, a trained model or a biological result.\par
\noindent Trainer and trainee name the current guidance roles. A machine-learning model can still be a learner; its fitted parameters do not measure a person's retained learning. Historical names and original source quotations remain historical.\par
\noindent Lookup labels help locate an entry. Some name related or contrasting concepts, so they are not interchangeable. The current JSON projection preserves every canonical term field, alias, source and definition history.\par
\noindent Physical scale and units, computational detail and statistical aggregation remain separate. Missing ranges or measurements stay unspecified. Both paper editions retain the same definitions and scientific scope; explanatory wording may differ.\par

% Stable term ID: a-checked-message-between-software-parts
\begin{samepage}
\subsection*{A checked message between software parts}
\noindent A structured message and interface specify what the values mean and how software parts exchange them.\par
\end{samepage}
\begin{samepage}
\noindent The checks include recording, ordered channels, quantity, timing where supplied, units and body axes. Passing a check does not prove that the caller’s claimed physical source is true.\par\nopagebreak[4]

\end{samepage}


\noindent\textit{Current registry lookup labels:} Translation packet, typed port and API; packet; translation port; typed interface; API; application programming interface\par

% Stable term ID: a-group-of-recorded-neural-features
\begin{samepage}
\subsection*{A group of recorded neural features}
\noindent Several recorded neural features considered together when predicting an outcome.\par
\end{samepage}
\begin{samepage}
\noindent Their joint pattern can help predict movement. That does not establish that the same pattern has the same meaning in another animal, task or day.\par\nopagebreak[4]

\end{samepage}


\noindent\textit{Current registry lookup labels:} Neural population; population decoding; population model\par

% Stable term ID: a-map-of-physical-connections
\begin{samepage}
\subsection*{A map of physical connections}
\noindent A structural map records physical connections. A functional recording measures activity over time, such as which recorded features produced events during a task.\par
\end{samepage}
\begin{samepage}
\noindent The structure can guide a model, but it cannot supply missing time-recorded activity. A drawing, connection map and synchronized functional recording have different evidential roles.\par\nopagebreak[4]

\end{samepage}


\noindent\textit{Current registry lookup labels:} Structural connectome and anatomical prior; connectome; connectomic connectivity; anatomical prior; structural connectivity; structural reference; recorded function; structural versus functional measurement\par

% Stable term ID: a-model-spanning-the-body
\begin{samepage}
\subsection*{A model spanning the body}
\noindent A future model connecting relevant observations and actions across the brain, spinal cord, nerves, muscles and movement.\par
\end{samepage}
\begin{samepage}
\noindent The project does not yet have a matched voluntary recording across that whole chain. Separate datasets cannot fill the missing measurements.\par\nopagebreak[4]

\end{samepage}


\noindent\textit{Current registry lookup labels:} Body-wide functional model; body-wide neural fusion; whole-body neural translator\par

% Stable term ID: a-specified-human-related-feature
\begin{samepage}
\subsection*{A specified human-related feature}
\noindent Studying a defined human-related biological component or feature in a model organism.\par
\end{samepage}
\begin{samepage}
\noindent This can overlap with adding donor cells. One selected immune or neural feature does not make an animal a miniature human or establish human-like cognition.\par\nopagebreak[4]

\end{samepage}


\noindent\textit{Current registry lookup labels:} Humanization; humanized; human-related feature\par

% Stable term ID: ability-that-remains-after-guidance-is-removed
\subsection*{Ability that remains after guidance is removed}
\noindent A learned ability that remains without the training guidance and can be tested again later.\par
\noindent Remove trainer commands while keeping any bridge needed to express the trainee’s own control. Switching off that required route can prevent movement regardless of what the trainee learned.\par

\noindent\textit{Current registry lookup labels:} Retained learning; retention; independent pupil control; retained skill; pupil-alone; necessary bridge held constant; pupil-alone with necessary bridge constant; independent learner control; learner-alone; learner-alone with necessary bridge constant; independent trainee control; trainee-alone; trainee-alone with necessary bridge constant\par

% Stable term ID: action2world
\begin{samepage}
\subsection*{Action2World}
\noindent The hand reaches; the cup moves or stays still. Contact, load and the measured outcome tell us whether the action did what we intended, and that feedback shapes the next attempt.\par
\end{samepage}
\noindent The relationship between an action and its consequences for the body and surrounding environment, including the sensory information those consequences produce.\par
\noindent A predicted consequence is not an observed outcome. Specify the object, reference frame, task, units and uncertainty. Footage supplies movement examples without supplying the performer's neural, muscle or force measurements.\par
\begin{samepage}
\noindent\textit{Organizational level:} whole system; network; signal/feature\par\nopagebreak[4]

\end{samepage}


\noindent\textit{Current registry lookup labels:} action2world; action to world; action-to-world\par

% Stable term ID: activity-state-inside-the-muscle-model
\begin{samepage}
\subsection*{Activity state inside the muscle model}
\noindent The model’s internal state that responds to a command and affects its force.\par
\end{samepage}
\begin{samepage}
\noindent This state can change more slowly than the command. The displayed activation comes from the computer model, not recovered physiological activation from the macaque recording.\par\nopagebreak[4]

\end{samepage}


\noindent\textit{Current registry lookup labels:} Muscle activation; activation; activation dynamics\par

% Stable term ID: adjustment-using-evidence-from-the-new-setting
\begin{samepage}
\subsection*{Adjustment using evidence from the new setting}
\noindent Changing or checking a model using examples from where it will be used.\par
\end{samepage}
\begin{samepage}
\noindent Adjusting a decoder, checking probability scores and measuring a body’s physical properties are different kinds of calibration. The new-day model uses labeled examples.\par\nopagebreak[4]

\end{samepage}


\noindent\textit{Current registry lookup labels:} Calibration; recalibration; calibrated transfer; uncalibrated; physiological calibration; physiological fidelity; subject calibration\par

% Stable term ID: agent2agent
\begin{samepage}
\subsection*{Agent2Agent}
\noindent One software agent passes declared information to another.\par
\end{samepage}
\noindent A proposed link carrying specified information between software agents.\par
\begin{samepage}
\noindent Name the agents, message or representation, direction and tested outcome. Software communication does not by itself supply a biological recording, sensory writer, shared experience or demonstrated skill transfer.\par\nopagebreak[4]

\end{samepage}


\noindent\textit{Current registry lookup labels:} agent2agent; agent-to-agent; agent to agent\par

% Stable term ID: agent2brain
\begin{samepage}
\subsection*{Agent2Brain}
\noindent An agent would give a person feedback during a stated task. Tests would have to distinguish help during that task from learning that remains afterwards.\par
\end{samepage}
\noindent A proposed path from a software agent to a person through specified feedback or a separately qualified neural output interface.\par
\begin{samepage}
\noindent Feedback would reach a person through an ordinary sensory route or a separately qualified direct neural interface. Those routes need different devices, evidence and admission: feedback on a screen is not direct neural delivery. EEG and EMG record; they do not write. The proposed agent-to-person path leaves the reverse observation and feedback direction intact. Help during a task must be tested separately from learning retained after it ends. Agent2Brain establishes no stimulation or physical-action authority. Physical spinal scope remains read-only.\par\nopagebreak[4]

\end{samepage}


\noindent\textit{Current registry lookup labels:} Agent2Brain?; agent2brain; agent2brain?; agent-to-brain; agent to brain\par

% Stable term ID: all-outputs-inside-their-ranges-together
\begin{samepage}
\subsection*{All outputs inside their ranges together}
\noindent How often every output of interest falls within its stated range in the same case.\par
\end{samepage}
\begin{samepage}
\noindent A 90\% range for each separate output does not guarantee 90\% coverage of all outputs at once or of a whole body.\par\nopagebreak[4]

\end{samepage}


\noindent\textit{Current registry lookup labels:} Simultaneous coverage; all-output coverage; whole-body coverage\par

% Stable term ID: an-independent-way-the-model-can-move
\begin{samepage}
\subsection*{An independent way the model can move}
\noindent One independent model movement, such as rotation at a joint. The formal shorthand is DOF.\par
\end{samepage}
\begin{samepage}
\noindent The native arm has seven coordinates and no separate finger movements. The human replay controls two goals while bracing other independent joints.\par\nopagebreak[4]

\end{samepage}


\noindent\textit{Current registry lookup labels:} Degree of freedom; DOF; DOFs; independent coordinate; braced joint\par

% Stable term ID: analysis-timing-and-biological-travel-time
\begin{samepage}
\subsection*{Analysis timing and biological travel time}
\noindent An analysis offset aligns two recorded quantities. A conduction delay concerns how long a signal travels through biology.\par
\end{samepage}
\begin{samepage}
\noindent The 40 ms sensory offset estimates an earlier limb state. The motor target’s 4 ms offset comes from sampling. Neither measures neural conduction time.\par\nopagebreak[4]

\end{samepage}


\noindent\textit{Current registry lookup labels:} S1 alignment offset and conduction delay; 40 ms sensory lag; S1 lag; 4 ms sampling offset; conduction delay; physiological latency\par

% Stable term ID: animal-drawing-and-measured-anatomy
\begin{samepage}
\subsection*{Animal drawing and measured anatomy}
\noindent An illustration explains a proposed animal or body. An empirical anatomical atlas is built from verified observations of anatomy; it cannot be established by a convincing drawing.\par
\end{samepage}
\begin{samepage}
\noindent The antechinus illustration is a proposal, separate from acquired macaque and human sources. The implementation has no matched validated antechinus neural/body atlas and makes no living-animal learning claim from the image.\par\nopagebreak[4]

\end{samepage}


\noindent\textit{Current registry lookup labels:} Antechinus proposal; antechinus; marsupial proposal; hypothetical organism; Proposed marsupial setting; anatomy illustration; illustrative model; empirical atlas; anatomical atlas\par

% Stable term ID: another-brain-area-used-in-the-finger-study
\begin{samepage}
\subsection*{Another brain area used in the finger study}
\noindent A brain area recorded alongside the movement area in the human finger study.\par
\end{samepage}
\begin{samepage}
\noindent The formal name is posterior parietal cortex, or PPC. The study provides attempted-finger labels, not measured trajectories of executed finger movement.\par\nopagebreak[4]

\end{samepage}


\noindent\textit{Current registry lookup labels:} Posterior parietal cortex; PPC; Parietal brain area\par

% Stable term ID: assumed-repair-to-the-native-muscle-model
\begin{samepage}
\subsection*{Assumed repair to the native muscle model}
\noindent A separately named model version replaces unusable zero-rate properties with the software’s default values.\par
\end{samepage}
\begin{samepage}
\noindent The source uses a named Schutte muscle model. Restoring 117 default properties permits numerical tests but does not measure the subject’s muscle response rates.\par\nopagebreak[4]

\end{samepage}


\noindent\textit{Current registry lookup labels:} Schutte muscle properties and default-rate fork; Schutte; zero-rate; assumed-rate fork; class defaults\par

% Stable term ID: attachment-order-and-the-chosen-route-segment
\begin{samepage}
\subsection*{Attachment order and the chosen route segment}
\noindent The ordered connection points and the numbered segment to which a wrapping rule applies.\par
\end{samepage}
\begin{samepage}
\noindent The rule uses the source model’s attachment order. The SDK counts from one and includes both range endpoints; the correction does not fit arbitrary points to an error frame.\par\nopagebreak[4]

\end{samepage}


\noindent\textit{Current registry lookup labels:} Attachment topology and SDK range; attachment topology; one-based inclusive range; humerus-to-ulna segment\par

% Stable term ID: axes-pushes-and-turns
\begin{samepage}
\subsection*{Axes, pushes and turns}
\noindent A coordinate frame states the axes used for movement. A handle wrench combines forces and turning moments measured in that frame.\par
\end{samepage}
\begin{samepage}
\noindent The handle record has three force and three moment components. An unknown mapping into the computer-body axes prevents treating those forces as verified simulator inputs.\par\nopagebreak[4]

\end{samepage}


\noindent\textit{Current registry lookup labels:} Coordinate frame and handle wrench; body frame; XY; physical frame; handle wrench; forces and moments\par

% Stable term ID: biological-tissue-repair
\begin{samepage}
\subsection*{Biological tissue repair}
\noindent Restoration or replacement of biological tissue or its function, supported by direct biological evidence.\par
\end{samepage}
\begin{samepage}
\noindent Better simulated tracking, assisted movement and retained behavior are different outcomes. The prototype does not demonstrate neural tissue repair.\par\nopagebreak[4]

\end{samepage}


\noindent\textit{Current registry lookup labels:} Regeneration; tissue regeneration; repair\par

% Stable term ID: biological-virtualization
\begin{samepage}
\subsection*{Biological virtualization}
\noindent For a movement model, matching the path is one question; matching the living system’s response to an intervention is another.\par
\end{samepage}
\noindent A claim that a digital model reproduces specified biological behaviour and causal dynamics within an independently tested domain and stated error bounds.\par
\noindent Independent ground truth must establish measured error and repeatability within declared conditions and tolerances. For causal dynamics, the comparison also needs independently measured responses to appropriate causal interventions; a similar-looking trajectory alone cannot establish them. Approximation can make a model plausible and may contribute to a validated model when its assumptions, errors and domain are tested and disclosed. A musculoskeletal simulation is itself a model to test, not independent empirical ground truth. Numerical detail and repeated verifier results guarantee neither biological accuracy nor complete physiology, subjective continuity, consciousness transfer or implemented Brain2Body capability. This biological claim is distinct from operating-system virtualization.\par
\noindent\textit{Scope:} validation: Specified observed behaviours and causal dynamical responses under declared conditions.\par
\noindent\textit{Organizational level:} whole system; network; organ; local circuit; cell; signal/feature\par
\noindent\textit{Computational granularity:} Declared model and brain-body interfaces\par
\begin{samepage}
\noindent\textit{Statistical aggregation:} Declared validation aggregation; separate from physical scale\par\nopagebreak[4]

\end{samepage}


\noindent\textit{Current registry lookup labels:} virtualization; virtualisation; biological virtualisation\par

% Stable term ID: body-dimensions-and-resistance-to-motion
\begin{samepage}
\subsection*{Body dimensions and resistance to motion}
\noindent Mass is the amount of matter modeled. Inertia describes resistance to changes in motion; anthropometry means body dimensions.\par
\end{samepage}
\begin{samepage}
\noindent The scaled model keeps assumed source masses and applies geometry-related inertia updates. Its length fit does not identify the animal’s actual masses, strengths or individual bone dimensions.\par\nopagebreak[4]

\end{samepage}


\noindent\textit{Current registry lookup labels:} Mass, inertia and anthropometry; inertia; mass; anthropometry; subject scaling; specimen scale\par

% Stable term ID: body-position-stored-by-the-simulator
\begin{samepage}
\subsection*{Body position stored by the simulator}
\noindent The model’s stored configuration, such as its joint positions.\par
\end{samepage}
\begin{samepage}
\noindent The formal variable name is qpos. Its entries depend on the model and are not automatically measured joint angles. The matching velocity variable, qvel, is explained separately.\par\nopagebreak[4]

\end{samepage}


\noindent\textit{Current registry lookup labels:} Generalized position and velocity; qpos; joint state\par

% Stable term ID: body-model-software-and-its-programming-tools
\begin{samepage}
\subsection*{Body-model software and its programming tools}
\noindent OpenSim provides tools for body models, geometry, scaling and simulation. Its software development kit, or SDK, lets code work with those models.\par
\end{samepage}
\begin{samepage}
\noindent Checks on these tools establish the stated model’s numerical behavior. They do not certify the exact upstream model version or the recorded animal’s anatomy.\par\nopagebreak[4]

\end{samepage}


\noindent\textit{Current registry lookup labels:} OpenSim and its SDK; OpenSim; SDK; software development kit; native SDK; OpenSim 4.5.2\par

% Stable term ID: body-position-feedback
\begin{samepage}
\subsection*{Body-position feedback}
\noindent Proprioception means sensing body position and movement. For example, it helps a body monitor where a limb is. Biofeedback uses bodily information to guide action or learning.\par
\end{samepage}
\begin{samepage}
\noindent The implemented correction uses simulated body position. It does not add measured peripheral nerve signals or show that a biological feedback process taught a retained skill.\par\nopagebreak[4]

\end{samepage}


\noindent\textit{Current registry lookup labels:} Proprioception and biofeedback; proprioception; simulated proprioception; biofeedback; Body-state sensing and feedback; proprioceptive feedback\par

% Stable term ID: body2action
\begin{samepage}
\subsection*{Body2Action}
\noindent The proposed link from body mechanics and sensory interaction to a specified, independently measured task action or outcome.\par
\end{samepage}
\noindent A task outcome makes this proposed link testable: the body’s forces, contacts and timing have to be measured in the relevant context. Body2Action connects body mechanics and sensory interaction to that specified action or outcome; it establishes no working controller. World labels are historical conceptual lookup references, not a strict synonym for every environmental interaction.\par
\noindent\textit{Organizational level:} Declared system or interface\par
\noindent\textit{Computational granularity:} Declared process/interface/model\par
\begin{samepage}
\noindent\textit{Statistical aggregation:} Validation domain stated independently of physical scale\par\nopagebreak[4]

\end{samepage}


\noindent\textit{Current registry lookup labels:} Body to action; Body2World; Body to world\par

% Stable term ID: body2body
% Canonical status: "user_requested_proposed_programme_role"
\subsection*{Body2Body}
\noindent A measured movement can constrain a digital body; a digital body can suggest a movement to test. Body2Body asks which quantities can travel between those settings and whether the resulting movement matches an independent reference.\par
\noindent A proposed research direction relating specified physical-body measurements to digital body models and testing mappings or feedback from those models to declared target bodies.\par
\begin{samepage}
\noindent The user’s physical-to-digital and digital-to-physical wording names a proposed direction, not a proven symmetric inverse. Observation-to-model fitting, retargeting and a separately qualified physical output have different requirements. Homologous and mechanically analogous features can support unpaired pretraining or transfer when their mechanics, scale, contacts, timing and missing features are declared. That learned relationship is not an actually measured neural pairing. Keep inferred and simulated quantities distinct from measurements. Biological virtualization still requires independent biological outcomes in a stated domain; no real body actuation, retained skill or whole-organism equivalence follows here.\par\nopagebreak[4]

\end{samepage}


\noindent\textit{Current registry lookup labels:} Body to body\par

% Stable term ID: brain-imaging-related-to-blood-oxygen
\begin{samepage}
\subsection*{Brain imaging related to blood oxygen}
\noindent An imaging method, called fMRI, that measures signals associated with changes in blood oxygenation.\par
\end{samepage}
\begin{samepage}
\noindent It does not record individual electrical spikes. The paper discusses a published model for this kind of image, not a project result showing that brain scans and spikes are interchangeable.\par\nopagebreak[4]

\end{samepage}


\noindent\textit{Current registry lookup labels:} Functional magnetic resonance imaging; fMRI\par

% Stable term ID: brain-pad
\subsection*{Brain pad}
\noindent A pad shows the proposed interface at the brain.\par
\noindent The proposed brain-side (cortical) interface, drawn as a pad.\par
\noindent A contextual display label in the user's Interoceptomics visual instructions, not a general clinical synonym. The name establishes no acquired neural data, registration, tissue access, function, safety, clinical use, species or physical dimensions.\par

\noindent\textit{Current registry lookup labels:} brain pad; cortical prosthesis; cortical interface; cortical interface prosthesis\par

% Stable term ID: brain-stimulation
\begin{samepage}
\subsection*{Brain stimulation}
\noindent An intervention that applies a specified physical stimulus to alter brain activity.\par
\end{samepage}
\noindent The method must be named, such as transcranial magnetic/electrical stimulation or intracortical microstimulation. These differ in access, risks and capabilities. EEG records; a scalp EEG headset does not provide general sensory writing or safe skill upload. Broad neuromodulation is not identical to precise task-content encoding.\par
\noindent\textit{Organizational level:} Declared game/model/interface role\par
\noindent\textit{Computational granularity:} Declared task-specific model or interface\par
\begin{samepage}
\noindent\textit{Statistical aggregation:} Validation scope declared per experiment; no physical-scale inference\par\nopagebreak[4]

\end{samepage}
\noindent\textit{Source:} A brain-computer interface that evokes tactile sensations improves robotic arm control\par
\noindent\textit{Source identifier:} flesher2021\par
\noindent\textit{DOI:} 10.1126/science.abd0380\par
\noindent\url{https://pubmed.ncbi.nlm.nih.gov/34016775/}\par
\noindent Research precedent for motor decoding plus sensory feedback through intracortical microstimulation; invasive and not a consumer EEG capability\par


\noindent\textit{Current registry lookup labels:} Neural stimulation; Neurostimulation\par

% Stable term ID: brain2agent
\begin{samepage}
\subsection*{Brain2Agent?}
\noindent A recording supplies observations that an agent may use for a stated task.\par
\end{samepage}
\noindent A proposed path from specified brain observations to a software agent through declared recording and interpretation interfaces.\par
\begin{samepage}
\noindent Recorded brain observations would reach an agent through a declared recording and interpretation interface. Where applicable, feedback toward the agent remains distinct from Agent2Brain’s proposed agent-to-person path. Neither direction identifies unique private intent or experience. Recording supplies observations; inference and simulation interpret or model them, while actual sensory delivery requires its own interface. Brain2Agent? remains an open proposal.\par\nopagebreak[4]

\end{samepage}


\noindent\textit{Current registry lookup labels:} Brain2Agent; brain2agent; brain2agent?; brain-to-agent; brain to agent\par

% Stable term ID: brain2body
\begin{samepage}
\subsection*{Brain2Body}
\noindent Connecting brain and body function; the 2 means to.\par
\end{samepage}
\noindent The numeral 2 is the preferred spelling. The programme includes body-to-brain feedback as well as outward pathways; the name does not establish complete biological translation or novelty of body–brain research.\par
\begin{samepage}
\noindent\textit{Organizational level:} whole system; network; local circuit; signal/feature\par\nopagebreak[4]

\end{samepage}


\noindent\textit{Current registry lookup labels:} brain2body; Brain2body; brain 2 body; brain-to-body; brain to body\par

% Stable term ID: brain2brain
\subsection*{Brain2Brain}
\noindent Information links between two brain or neural systems.\par
\noindent The project’s label for studying specified information links between brain or neural systems, including trainer–trainee comparisons.\par
\noindent Brain2Brain is the numeral-style public label for the project’s existing brain-to-brain research direction. Each experiment must specify its signal, task and controls; the name itself establishes no shared thoughts, general skill transfer or successful biological teaching.\par
\noindent\textit{Organizational level:} whole system; network; local circuit; signal/feature\par

\noindent\textit{Current registry lookup labels:} brain2brain; brain 2 brain; brain-to-brain; brain to brain\par

% Stable term ID: brain2motion
\begin{samepage}
\subsection*{Brain2Motion}
\noindent The proposed movement pathway connecting Brain2Body with Body2Action and the resulting measured motion, including the returning sensory feedback.\par
\end{samepage}
\noindent The label concerns the movement part of the wider programme, not every sensory/cognitive function. Brain/spinal/peripheral observations, muscle activity, body mechanics and external outcomes stay distinct. A decoded motion signal, simulation or lifelike animation is not a complete physiological explanation or validated skill transfer.\par
\noindent\textit{Organizational level:} Declared system or interface\par
\noindent\textit{Computational granularity:} Declared process/interface/model\par
\begin{samepage}
\noindent\textit{Statistical aggregation:} Validation domain stated independently of physical scale\par\nopagebreak[4]

\end{samepage}


\noindent\textit{Current registry lookup labels:} Brain to motion\par

% Stable term ID: brain2spine
\begin{samepage}
\subsection*{Brain2Spine}
\noindent A movement command meets a spinal cord that already regulates reflexes, coordinates muscles and adapts through experience. Understanding that middle step helps explain how the same intention can produce different actions.\par
\end{samepage}
\noindent The pathway from brain activity into spinal processing: descending signals influence circuits that combine them with sensory input and the body's current state.\par
\noindent A scalp EEG recording does not identify every descending signal or measure the spinal circuit's hidden state. The physical spine remains read only. This label does not name an admitted stimulation interface.\par
\begin{samepage}
\noindent\textit{Organizational level:} whole system; network; signal/feature\par\nopagebreak[4]

\end{samepage}


\noindent\textit{Current registry lookup labels:} brain2spine; brain to spine; brain-to-spine\par

% Stable term ID: brainnet
\begin{samepage}
\subsection*{BrainNet}
\noindent A network of specified links between brain or neural systems.\par
\end{samepage}
\noindent A working network metaphor for connected brain or neural systems within the programme.\par
\noindent BrainNet pictures specified links between brain or neural systems as a network. It is a programme metaphor; it establishes no validated universal neural network, uniquely coined field or available trademark.\par
\begin{samepage}
\noindent\textit{Organizational level:} whole system; network; local circuit; signal/feature\par\nopagebreak[4]

\end{samepage}


\noindent\textit{Current registry lookup labels:} brainnet; Brain Net\par

% Stable term ID: brain-body-interface
% Canonical status: "user_requested_proposed_programme_role"
\subsection*{Brain–Body Interface}
\noindent A brain recording and a measured hand movement answer different questions about the same reach. A Brain–Body Interface would connect the stated observations and any separately qualified outputs so that the proposed relationship can be tested.\par
\noindent A proposed interface linking specified brain observations with body observations or control outputs in a declared task, with each direction and physical route identified.\par
\begin{samepage}
\noindent This names an interface role, rather than the Brain2Body programme or a model spanning the body. A Digital Bridge is the existing selected-information route around an interruption or between specified recording and control channels; it is related, not an automatic synonym for every Brain–Body Interface. Recording, inference, simulated control and physical output require distinct evidence. EEG and EMG record; a feedback or stimulation output needs its own qualified interface. The label establishes neither complete coverage nor a bidirectional biological loop. No new acronym is introduced.\par\nopagebreak[4]

\end{samepage}


\noindent\textit{Current registry lookup labels:} Brain-Body Interface\par

% Stable term ID: choosing-a-model-and-deciding-if-it-passes
\begin{samepage}
\subsection*{Choosing a model and deciding if it passes}
\noindent Choosing settings and applying a stated pass rule using the validation data.\par
\end{samepage}
\begin{samepage}
\noindent The frozen-transfer rule requires a prediction score at least 0.01 above every specified control. A failing candidate stays rejected rather than being given an easier success rule.\par\nopagebreak[4]

\end{samepage}


\noindent\textit{Current registry lookup labels:} Validation selection and admission; validation-selected; acceptance criterion; admission; failure criterion\par

% Stable term ID: command-sent-to-a-modeled-muscle
\begin{samepage}
\subsection*{Command sent to a modeled muscle}
\noindent The input command that starts the modeled muscle’s activation process.\par
\end{samepage}
\begin{samepage}
\noindent The command is bounded and nonnegative here. It differs from the resulting activity state, recorded muscle electricity and instructions for a biological intervention.\par\nopagebreak[4]

\end{samepage}


\noindent\textit{Current registry lookup labels:} Muscle excitation; excitation; actuator command; muscle control\par

% Stable term ID: communication-route-around-an-interruption
\begin{samepage}
\subsection*{Communication route around an interruption}
\noindent An interface carrying selected information around a modeled break or between stated recording and control channels.\par
\end{samepage}
\begin{samepage}
\noindent A bridge can compensate without repairing tissue or supplying muscle energy. The virtual interruption does not establish a biological injury model or treatment.\par\nopagebreak[4]

\end{samepage}


\noindent\textit{Current registry lookup labels:} Digital bridge; compensatory bridge; brain–spinal bridge; BCI; brain-computer interface; virtual pathway interruption\par

% Stable term ID: comparable-roles-in-different-organisms
\begin{samepage}
\subsection*{Comparable roles in different organisms}
\noindent Comparing roles such as reaching or support while keeping each organism’s anatomy explicit.\par
\end{samepage}
\begin{samepage}
\noindent Related roles do not imply identical neurons or muscles. A combined fictional body cannot replace the measured sources and checks needed for each real organism.\par\nopagebreak[4]

\end{samepage}


\noindent\textit{Current registry lookup labels:} Organism role and homology; organism adapter; limb role; functional coordinate; homology; biped/quadruped\par

% Stable term ID: comparisons-that-test-whether-guidance-matters
\subsection*{Comparisons that test whether guidance matters}
\noindent Matched replay uses the relevant recorded activity under matched task conditions. Stale or unmatched replay uses activity from a different time or condition. No-link and inactive or sham comparisons test whether the proposed information route adds benefit.\par
\noindent These are proposed comparison types, not new animal results. Shuffling disrupts timing; a yoked comparison gives matched help from another case. Guidance amount, timing and later trainee control need their own stated comparisons.\par

\noindent\textit{Current registry lookup labels:} Shuffled and yoked controls; shuffled-input control; time-shuffled; yoked signal comparison; Matched comparisons with altered information; matched replay; matched own-session activity; stale replay; unmatched replay; no-link; sham; inactive comparison\par

% Stable term ID: completed-time-interval-for-counting-events
\begin{samepage}
\subsection*{Completed time interval for counting events}
\noindent A stated time interval groups events into one count. A half-open interval includes its start and excludes its end.\par
\end{samepage}
\begin{samepage}
\noindent Inputs use completed intervals. The exact boundary rule, interval duration and timing must remain explicit rather than being guessed from the display.\par\nopagebreak[4]

\end{samepage}


\noindent\textit{Current registry lookup labels:} Completed time bin and half-open binning; bin; 20 ms bin; 50 ms bin; half-open; bin end; cadence\par

% Stable term ID: computer-processors-memory-and-file-size
\begin{samepage}
\subsection*{Computer processors, memory and file size}
\noindent A CPU is a general processor, a GPU is another kind of parallel processor, and RAM is working memory. Bytes and GiB measure digital data size.\par
\end{samepage}
\begin{samepage}
\noindent The training described in this source edition used four CPU threads; GPU acceleration was neither verified nor used. AI-assisted coding does not establish independent biological validation or an authorship credential.\par\nopagebreak[4]

\end{samepage}


\noindent\textit{Current registry lookup labels:} Computing hardware and data size; CPU; GPU; RAM; AMD; AI; Ryzen AI Max+395; GiB; byte; bytes\par

% Stable term ID: correct-label-rate-and-chance-comparison
\begin{samepage}
\subsection*{Correct-label rate and chance comparison}
\noindent Accuracy is the proportion of labels chosen correctly. A confusion table shows which labels were mistaken for others.\par
\end{samepage}
\begin{samepage}
\noindent A uniform guess among ten balanced classes has a 10\% reference rate. Accuracy and label scores do not prove executed movement or confidence in biological truth.\par\nopagebreak[4]

\end{samepage}


\noindent\textit{Current registry lookup labels:} Classification accuracy and chance; accuracy; uniform chance; ten-class confusion; confusion matrix\par

% Stable term ID: count-prediction-improvement-per-recorded-event
\begin{samepage}
\subsection*{Count-prediction improvement per recorded event}
\noindent Bits per spike describes how much a count-prediction score improves over a defined baseline, divided by the number of recorded spikes.\par
\end{samepage}
\begin{samepage}
\noindent The baseline uses the training mean rate. A positive result concerns estimated average counts, not newly recorded activity or verified stimulation instructions.\par\nopagebreak[4]

\end{samepage}


\noindent\textit{Current registry lookup labels:} Bits per spike; BPS; ΔBPS; bits/spike; likelihood gain\par

% Stable term ID: data-for-learning-choosing-and-checking
\begin{samepage}
\subsection*{Data for learning, choosing and checking}
\noindent Separate portions used to fit the model, choose its settings, and check the selected result.\par
\end{samepage}
\begin{samepage}
\noindent The paper uses local time-ordered splits. If test results have already been inspected, later reuse cannot be described as another independent test.\par\nopagebreak[4]

\end{samepage}


{\rightskip=0pt plus1fil\relax
\noindent\textit{Current registry lookup labels:} Training, validation and test partitions; train/validation/test; train/val/test; split; chronological 70/15/15\par
}

% Stable term ID: data-kept-out-of-a-stated-step
\begin{samepage}
\subsection*{Data kept out of a stated step}
\noindent Data reserved from training or from another specified development step.\par
\end{samepage}
\begin{samepage}
\noindent Independence depends on what was actually excluded. A previously inspected case does not become a fresh test just because it is called held-out.\par\nopagebreak[4]

\end{samepage}


\noindent\textit{Current registry lookup labels:} Holdout; held-out; reserved window; independent evaluation\par

% Stable term ID: describing-motion-and-calculating-what-causes-it
\begin{samepage}
\subsection*{Describing motion and calculating what causes it}
\noindent Describing a body’s position and movement is one task; calculating how modeled forces change that movement is another.\par
\end{samepage}
\begin{samepage}
\noindent The paper calls these kinematics and dynamics. A replayed pose and a muscle-driven trajectory are different results; neither alone proves physiological accuracy.\par\nopagebreak[4]

\end{samepage}


\noindent\textit{Current registry lookup labels:} Kinematics and dynamics; kinematics; dynamics; forward dynamics; numerical plant\par

% Stable term ID: ecotype
% Canonical status: "established_scientific_reference"
\subsection*{Ecotype}
\noindent A population adapted to a coastal habitat is an example of the ecological distinction at issue. Ecotype names the adapted population or form; the surroundings themselves are its environment.\par
\noindent A locally adapted population or form associated with a particular ecological setting.\par
\noindent Ecotype is an established population/ecological concept, not an individual’s surroundings or an obligatory third stage after genotype and phenotype. State the organism and evidence for the claimed local differentiation or adaptation. The retained Turesson source is a historical plant example~\citep{turesson1922ecotype}, not a complete modern evolutionary theory or evidence about the proposed nervous-system experiments.\par


\noindent\textit{Current registry lookup labels:} No additional label declared.\par

% Stable term ID: electrical-recording-in-outer-brain-tissue
\begin{samepage}
\subsection*{Electrical recording in outer brain tissue}
\noindent Electrodes in the brain’s outer layer, called the cortex, record electrical signals outside the cells.\par
\end{samepage}
\begin{samepage}
\noindent This differs from a scalp recording. It also does not observe every cell in a circuit; the source file defines which events and channels were recorded.\par\nopagebreak[4]

\end{samepage}


\noindent\textit{Current registry lookup labels:} Intracortical extracellular recording; intracortical; extracellular cortical events; Electrical recording inside the cortex; cortex\par

% Stable term ID: enactome
% Canonical status: "user_requested_proposed_data_organization"
\subsection*{Enactome}
\noindent A reach has a path, a time and a result. An Enactome would organize observations like these so that the action and what followed can be compared across stated conditions.\par
\noindent A proposed organization of specified action observations and the recorded relationships between actions and their outcomes in declared tasks.\par
\noindent Declare which actions, body parts, tasks, units, frames and times the data cover. Keep commands, measured movement, sensory consequences and simulated outcomes separate. An observed association does not establish causation, a complete repertoire or the person’s unique intent. This proposed data organization differs from the tentative Enactotype individual-profile label.\par


\noindent\textit{Current registry lookup labels:} No additional label declared.\par

% Stable term ID: endohumanisation
% Canonical status: "user_requested_proposed_programme_role"
\subsection*{EndoHumanisation}
\noindent Choose a particular human-related feature, then ask what introducing or studying that feature would change in the model organism. EndoHumanisation names that proposed line of work.\par
\noindent A proposed research stage studying specified human-related genetic features within a model organism and testing their relevance to a declared biological question.\par
\noindent The roadmap places XenoTransgenesis at step 1.1. The working interpretation concerns defined genetic features; the actual manipulation and its established scientific name remain to be specified. A selected feature does not make an animal human or establish human cognition, disease fidelity or an appropriate intervention age. Biological relevance, function and animal/participant, ethical, safety and rights requirements need their own evidence and admission. This entry supplies no experimental procedure.\par


\noindent\textit{Current registry lookup labels:} No additional label declared.\par

% Stable term ID: error-range-based-on-validation-examples
\begin{samepage}
\subsection*{Error range based on validation examples}
\noindent A range placed around a prediction using the errors seen on the data used for model selection.\par
\end{samepage}
\begin{samepage}
\noindent The ranges have constant width for each output. They do not guarantee future coverage when examples overlap in time or their errors are related.\par\nopagebreak[4]

\end{samepage}


\noindent\textit{Current registry lookup labels:} Residual band; empirical residual bounds; validation residual quantile; 90\% per-output residual band\par

% Stable term ID: events-grouped-as-one-neural-feature
\begin{samepage}
\subsection*{Events grouped as one neural feature}
\noindent The source recording process groups events into features that the model can count separately.\par
\end{samepage}
\begin{samepage}
\noindent A group is not a verified lifelong identity for one neuron. It must stay separate from unassigned events and from the channel labels used in other recordings.\par\nopagebreak[4]

\end{samepage}


\noindent\textit{Current registry lookup labels:} Sorted neural unit; sorted unit; sorted cell; sorted feature\par

% Stable term ID: exohumanisation
% Canonical status: "user_requested_proposed_programme_role"
\subsection*{ExoHumanisation}
\noindent Added tissue can survive without doing the job expected of it. ExoHumanisation names a proposed study of a specified human-related graft: what survives, what connects, and what function can actually be measured?\par
\noindent A proposed research stage studying specified human-related donor cells or tissue in a model organism and testing their survival, integration and function.\par
\noindent The roadmap associates step 2.1 with XenoTransplantation. A cross-species graft retains that established meaning; cell survival, anatomical connection, activity and useful functional integration are distinct outcomes. A graft does not establish a human nervous system, a shared mind or whole-organism human equivalence. Species, tissue identity, biological purpose, ethics, safety and rights remain explicit dependencies. This label provides no transplantation procedure or intervention-age window.\par


\noindent\textit{Current registry lookup labels:} No additional label declared.\par

% Stable term ID: expected-event-count-and-its-fit-score
\begin{samepage}
\subsection*{Expected event count and its fit score}
\noindent A model estimates the average event count for given inputs, then a score checks how that estimate fits the recorded counts.\par
\end{samepage}
\begin{samepage}
\noindent The paper uses Poisson negative log likelihood for that check. A predicted average may be fractional; a good count-fit score does not establish a complete biological spike model.\par\nopagebreak[4]

\end{samepage}


\noindent\textit{Current registry lookup labels:} Conditional count mean and Poisson likelihood; expected count mean; conditional count; Poisson; negative log likelihood; NLL\par

% Stable term ID: external-movement-model
\begin{samepage}
\subsection*{External movement model}
\noindent The external part of the body model, constrained by independently measured surface movement and pose in a declared task and coordinate frame.\par
\end{samepage}
\noindent The motion-capture relay-composite retains the source views, clocks, calibration and coordinate frame for the declared task. The earlier upper/lower hand concept and the later synchronized left/right tripod arm proposal are different arrangements; each needs its own visibility, synchronization and error checks. The composite is derived. Hidden internal states, forces and intent are not uniquely measured. LiDAR and optical infrared tracking require distinct hardware, calibration and uncertainty qualification.\par


\noindent\textit{Current registry lookup labels:} No additional label declared.\par

% Stable term ID: extra-correction-and-its-scoring-rule
\begin{samepage}
\subsection*{Extra correction and its scoring rule}
\noindent A model can add a correction to an existing controller. A reward scores the result; command-change rate describes how quickly the input changes.\par
\end{samepage}
\begin{samepage}
\noindent The formal names include residual policy, reinforcement learning and actuator slew. A simulator can score tracking or task completion without proving biological learning or correct hidden neural activity.\par\nopagebreak[4]

\end{samepage}


\noindent\textit{Current registry lookup labels:} Residual policy, reward and actuator slew; residual policy; reinforcement learning; RL; verifiable reward; actuator slew; task completion\par

% Stable term ID: file-fingerprints
\begin{samepage}
\subsection*{File fingerprints}
\noindent A checksum or hash is calculated from file bytes to detect changes or match a stated reference.\par
\end{samepage}
\begin{samepage}
\noindent The named methods here include MD5 and SHA-256. A match checks bytes against that reference; it does not prove scientific truth or verify an entire file when only a range was acquired.\par\nopagebreak[4]

\end{samepage}


\noindent\textit{Current registry lookup labels:} Checksums and hashes; checksum; hash; MD5; SHA-256; source checksum; archive hash\par

% Stable term ID: force-in-the-modeled-muscle-and-tendon
\begin{samepage}
\subsection*{Force in the modeled muscle and tendon}
\noindent The force produced by the computer model for a muscle-tendon unit, or a force measured by a stated physical method.\par
\end{samepage}
\begin{samepage}
\noindent The native test reads forces from the running simulator rather than an old lookup table. Those values are still simulated rather than isolated forces recorded in the source animal.\par\nopagebreak[4]

\end{samepage}


\noindent\textit{Current registry lookup labels:} Tendon force; muscle force; isolated muscle force\par

% Stable term ID: form-function-action
\begin{samepage}
\subsection*{Form \(\rightarrow\) function \(\rightarrow\) action}
\noindent The programme narrative follows structural constraints, the context-dependent processes they support, and independently measured actions or outcomes.\par
\end{samepage}
\noindent The three words organize the explanation; every molecular and neural level has its own structure and dynamics. Form informs and constrains function but does not determine the full process alone. Feedback returns outcomes to the system; this is not a one-way causal chain or a guaranteed one-to-one mapping.\par
\noindent\textit{Organizational level:} Declared system or interface\par
\noindent\textit{Computational granularity:} Declared process/interface/model\par
\begin{samepage}
\noindent\textit{Statistical aggregation:} Validation domain stated independently of physical scale\par\nopagebreak[4]

\end{samepage}


\noindent\textit{Current registry lookup labels:} Form to function to action\par

% Stable term ID: functional-electrical-stimulation
\begin{samepage}
\subsection*{Functional electrical stimulation}
\noindent Electrical stimulation coordinated to produce a specified functional movement or task.\par
\end{samepage}
\noindent FES is the appropriate functional-actuation label for the proposed later assisted-body study. It requires separately qualified devices, control/safety limits and specialist-supervised research approval. This registry contains no stimulation settings or self-experiment instructions.\par
\noindent\textit{Organizational level:} Declared game/model/interface role\par
\noindent\textit{Computational granularity:} Declared task-specific model or interface\par
\begin{samepage}
\noindent\textit{Statistical aggregation:} Validation scope declared per experiment; no physical-scale inference\par\nopagebreak[4]

\end{samepage}
\noindent\textit{Source:} Restoring cortical control of functional movement in a human with quadriplegia\par
\noindent\textit{Source identifier:} bouton2016\par
\noindent\textit{DOI:} 10.1038/nature17435\par
\noindent\url{https://www.nature.com/articles/nature17435}\par
\noindent Research precedent linking implanted cortical recording to a custom neuromuscular stimulation sleeve; not an EEG/OpenBCI or elite-skill result\par


\noindent\textit{Current registry lookup labels:} FES\par

% Stable term ID: further-training-for-the-target-task
\begin{samepage}
\subsection*{Further training for the target task}
\noindent Continuing model training so it predicts the desired target, using the relevant inputs and labels.\par
\end{samepage}
\begin{samepage}
\noindent The paper calls this fine-tuning. Changing software weights is not evidence that an animal or person learned a biological skill.\par\nopagebreak[4]

\end{samepage}


\noindent\textit{Current registry lookup labels:} Fine-tuning; paired-output fine-tuning\par

% Stable term ID: future-testing-and-analysis-of-existing-data
\begin{samepage}
\subsection*{Future testing and analysis of existing data}
\noindent A future test checks later unseen outcomes under a fixed plan; an existing-data test analyzes outcomes already recorded.\par
\end{samepage}
\begin{samepage}
\noindent The second source file was chosen before inspection, but its recording day is earlier than the original model’s day. Choosing a file in advance does not make it future biological deployment.\par\nopagebreak[4]

\end{samepage}


\noindent\textit{Current registry lookup labels:} Prospective and retrospective evaluation; prospective; retrospective; prospective biological deployment\par

% Stable term ID: gap-between-data-portions
\begin{samepage}
\subsection*{Gap between data portions}
\noindent Leaving out an interval around a split boundary so overlapping windows do not cross into another portion.\par
\end{samepage}
\begin{samepage}
\noindent The continuous pipeline uses a one-second gap and checks the input and target support. This protects the software split, not every assumption about the original measurement.\par\nopagebreak[4]

\end{samepage}


\noindent\textit{Current registry lookup labels:} Partition embargo; embargo; boundary embargo\par

% Stable term ID: genotype
% Canonical status: "established_scientific_reference"
\subsection*{Genotype}
\noindent A genetic description can record which versions of a DNA sequence an organism carries. Genotype names that genetic make-up, whether the description covers a selected part of the DNA or a wider set. The scope must be stated; the genetic description alone does not tell us every observable trait.\par
\noindent The genetic constitution being described for an organism, with the genetic scope of the description stated.\par
\noindent Genotype is an established scientific term. Its description is not a complete prediction of the organism’s observed traits; phenotype also depends on environment and development. The three labels in the user’s organism-level example are not three obligatory sequential stages or three fixed individual types.\par


\noindent\textit{Current registry lookup labels:} No additional label declared.\par

% Stable term ID: hallucination
\begin{samepage}
\subsection*{Hallucination}
\noindent A model can predict a movement correctly while inventing the biological mechanism it claims produced it.\par
\end{samepage}
\noindent Invented or ungrounded model content presented as if it corresponded to reality. In programme claims, this includes fabricated biological structure, dynamics or ability presented as established biological fidelity.\par
\noindent Research on AlphaFold3 uses the term for generative-model errors; this programme also applies it to unsupported biological attributions. A disclosed approximation, hypothesis, counterfactual or unvalidated candidate is a different kind of statement, and missing evidence alone does not prove it false. Numerical error concerns the calculation; objective misspecification concerns what is being optimized; causal attribution error concerns the mechanism being credited. Those differences matter when deciding what a result supports. An ungrounded generated answer can happen to be correct without validating its claimed biological mechanism or reliability elsewhere. Informed predictions therefore still require independent performance evaluation. The rolling-dice analogy describes reliance on unverified claims, not literal randomness or equal odds. Calling every approximation a hallucination, or every connectome or simulation fictitious, would erase the distinction.\par
\noindent\textit{Scope:} usage: Model-generated biological content and claims of biological fidelity.\par
\begin{samepage}
\noindent\textit{Computational granularity:} Model output and its interpretation\par\nopagebreak[4]

\end{samepage}


\noindent\textit{Current registry lookup labels:} hallucination; model hallucination\par

% Stable term ID: help-during-current-performance
\subsection*{Help during current performance}
\noindent Help from a controller, replay, trainer or interface that improves performance while it is present.\par
\noindent For example, a trainee can move successfully while following trainer commands. This does not show independent learning or tissue regeneration.\par

\noindent\textit{Current registry lookup labels:} Assistance; compensation; assisted performance\par

% Stable term ID: how-a-muscle-s-pull-turns-a-joint
\begin{samepage}
\subsection*{How a muscle’s pull turns a joint}
\noindent To relate a muscle’s pulling force to joint turning, the model uses its geometry. Consistent work calculations also relate those forces to small coordinate changes.\par
\end{samepage}
\begin{samepage}
\noindent The formal terms are moment arm and virtual work. Their numerical checks do not measure biological muscle strength; the human controller uses them to account for dependent shoulder movements.\par\nopagebreak[4]

\end{samepage}


\noindent\textit{Current registry lookup labels:} Moment arm and virtual work; moment arm; virtual-work projection; dependent shoulder-girdle load\par

% Stable term ID: how-many-labeled-examples-were-supplied
\begin{samepage}
\subsection*{How many labeled examples were supplied}
\noindent The labeled examples given to train or choose a model, including examples from prior recordings.\par
\end{samepage}
\begin{samepage}
\noindent The first 5\% counts new training labels only. Validation labels, earlier-source labels and extra neural samples are additional information.\par\nopagebreak[4]

\end{samepage}


\noindent\textit{Current registry lookup labels:} Label budget and supervision; label budget; paired supervision; adaptation labels; validation labels; source-task supervision\par

% Stable term ID: how-scores-from-several-outputs-are-combined
\begin{samepage}
\subsection*{How scores from several outputs are combined}
\noindent Scores can be averaged across individual outputs or first combined within defined groups. In this entry, macro describes the averaging and weighting rule, not the size of a body part, the scope of a nervous system, or the level of detail in a computer model.\par
\end{samepage}
\begin{samepage}
\noindent The sensory result gives equal weight to five physical output groups. This avoids letting a larger group dominate but does not make their different physical units interchangeable.\par\nopagebreak[4]

\end{samepage}


\noindent\textit{Current registry lookup labels:} Per-output and group-macro scoring; per-output; group-macro; macro R²\par

% Stable term ID: how-uncertain-the-result-remains
\begin{samepage}
\subsection*{How uncertain the result remains}
\noindent The remaining doubt about a prediction, including why it could be wrong or change.\par
\end{samepage}
\begin{samepage}
\noindent Noise in the instrument, errors in the model, changing conditions and missing observations are different causes. The current ranges do not separate all of them.\par\nopagebreak[4]

\end{samepage}


\noindent\textit{Current registry lookup labels:} Prediction uncertainty; uncertainty; model uncertainty; measurement noise; drift\par

% Stable term ID: human-arm-model-used-for-playback
\begin{samepage}
\subsection*{Human arm model used for playback}
\noindent MyoArm is the human arm model used for the displayed muscle-driven replay.\par
\end{samepage}
\begin{samepage}
\noindent The replay uses 63 modeled muscles and two mapped joint goals. Those quantities are simulated, and the macaque velocity mapping does not prove neural or skill transfer into a human body.\par\nopagebreak[4]

\end{samepage}


\noindent\textit{Current registry lookup labels:} MyoArm; human MyoArm; MyoArm right\par

% Stable term ID: internal-musculoskeletal-model
\begin{samepage}
\subsection*{Internal musculoskeletal model}
\noindent The internal anatomy part of the body model: specified bones, joints and soft tissues, assembled from appropriate registered anatomical measurements.\par
\end{samepage}
\noindent Proposed complementary X-ray/CT and MRI inform different structures. Anatomy does not directly measure activation, force, material properties or neural function. Wētā FX provides visual/simulation inspiration, not a verified identical acquisition pipeline.\par


\noindent\textit{Current registry lookup labels:} No additional label declared.\par

% Stable term ID: interoception
\begin{samepage}
\subsection*{Interoception}
\noindent Information about the body’s internal condition, such as signals related to heartbeat or breathing.\par
\end{samepage}
\begin{samepage}
\noindent\textit{Scope:} Sensing the body’s internal condition differs from sensing a limb’s position and movement. Internal sensing remains a research objective; this prototype does not provide a fully measured body-wide internal signal.\par\nopagebreak[4]

\end{samepage}


\noindent\textit{Current registry lookup labels:} Sensing the body from inside; interoception; interoceptive signal; internal sensations\par

% Stable term ID: interoceptome
% Canonical status: "user_requested_proposed_data_organization"
\subsection*{Interoceptome}
\noindent A heartbeat signal and a breathing signal describe different parts of the body’s internal condition. An Interoceptome would organize the selected observations and their relationships, retaining how, where and when each was recorded.\par
\noindent A proposed organization of specified observations of internal-body sensing and the relationships being studied among them.\par
\noindent This is a data-organization proposal within Interoceptomics. Its coverage must name the organism, body systems, measurements, timing and task. Measured signals, calculated quantities and model estimates remain distinguishable. Touch and limb-position sensing are not relabelled as interoception, and the suffix does not establish a complete internal-body map or an individual biological type.\par


\noindent\textit{Current registry lookup labels:} No additional label declared.\par

% Stable term ID: interoceptomics
\begin{samepage}
\subsection*{Interoceptomics}
\noindent How does the body sense a heartbeat or its breathing, and how are the nerve connections involved arranged? Interoceptomics brings those questions together in a proposed study, relating observations of internal-body sensing to maps of the connections carrying or processing that information.\par
\end{samepage}
\noindent A proposed field combining interoception and connectomics to study internal-body sensing and its relationships with neural connectivity across scales.\par
\noindent The programme uses Interoceptomics as an umbrella for this proposed research. Giving the work a name does not establish a validated comprehensive map or a recognized external standard.\par
\noindent\textit{Scope:} macro: Whole-body, organ and network relationships relevant to internal-body sensing; micro: Local signals, cells, circuits and interfaces relevant to those relationships.\par
\begin{samepage}
\noindent\textit{Organizational level:} whole system; network; organ; local circuit; cell; signal/feature\par\nopagebreak[4]
\nopagebreak[4]
\end{samepage}


\noindent\textit{Current registry lookup labels:} interoceptomics\par

% Stable term ID: keeping-each-source-separate
\begin{samepage}
\subsection*{Keeping each source separate}
\noindent Keeping track of which recording, channel, time and physical setting a value belongs to.\par
\end{samepage}
\begin{samepage}
\noindent A channel numbered 12 in one day is not automatically the same cell as channel 12 on another day. A software label also cannot prove where supplied data physically came from.\par\nopagebreak[4]

\end{samepage}


\noindent\textit{Current registry lookup labels:} Source conditioning; source-conditioned; source-specific; recording-specific\par

% Stable term ID: keeping-each-trial-together
\begin{samepage}
\subsection*{Keeping each trial together}
\noindent Assigning an entire trial to one data portion instead of dividing it among portions.\par
\end{samepage}
\begin{samepage}
\noindent This prevents pieces of the same trial from being shared between the relevant splits. It does not remove prior inspection or other source limitations.\par\nopagebreak[4]

\end{samepage}


\noindent\textit{Current registry lookup labels:} Whole-trial separation; whole-trial-separated; successful trial; trial partition\par

% Stable term ID: keeping-up-now-and-computing-for-later
\begin{samepage}
\subsection*{Keeping up now and computing for later}
\noindent Real-time computation keeps pace with incoming data; offline computation can take longer and be played back later.\par
\end{samepage}
\begin{samepage}
\noindent Native planning takes about 4.52 seconds for a 20 ms simulated step and is offline. Smooth cached playback does not prove live planning or biological control.\par\nopagebreak[4]

\end{samepage}


\noindent\textit{Current registry lookup labels:} Real-time and offline execution; real-time; offline; paced playback; cached reference\par

% Stable term ID: kind-of-measurement
\begin{samepage}
\subsection*{Kind of measurement}
\noindent What a method records, such as electrical spikes, scalp voltage, muscle electrical activity or a brain-scan signal.\par
\end{samepage}
\begin{samepage}
\noindent These signals measure different things. They cannot be treated as interchangeable numbers merely because all of them relate to the nervous system.\par\nopagebreak[4]

\end{samepage}


\noindent\textit{Current registry lookup labels:} Signal modality; modality; observation process\par

% Stable term ID: large-movement-model
\subsection*{Large Movement Model}
\noindent The shorthand ``NeuroTranslate with MoCap MyoSim = Large Movement Model'' names the proposed combination of NeuroTranslate, MoCap ReComp and MyoSim Replay. It brings a learned mapping, measured movement and virtual mechanics into one architecture. The equals sign is not a mathematical identity or evidence of a trained model. EEG, EMG, measured movement, virtual mechanics and the trainee retain their distinct roles.\par
\noindent The proposed movement-focused learned model that brings neural structure and functional observations together with internal musculoskeletal anatomy and independently measured external movement in a declared task.\par
\noindent The programme model-family name, analogous in ambition to large language model. It does not replace accurate technical architecture descriptions or mean that every DNN is a movement model. Scans and capture alone do not train a model. Parameter count, scale, generalization and nervous-system fidelity require evidence. System 3 thinking remains the partnership framing. In the October 19 shorthand, MoCap MyoSim combines MoCap ReComp and MyoSim Replay with NeuroTranslate. The equals sign names the proposed architecture; it is not a mathematical identity or evidence that a trained model already exists. These three named groups do not remove the distinct EEG, EMG, measured-movement, virtual-mechanics and trainee roles.\par

\noindent\textit{Current registry lookup labels:} large movement models; LMM; Large movement model\par

% Stable term ID: lasting-change-in-the-nervous-system
\begin{samepage}
\subsection*{Lasting change in the nervous system}
\noindent Changes in the nervous system that can contribute to learning or adaptation.\par
\end{samepage}
\begin{samepage}
\noindent Updating software weights or controller gains is a different event. The practice interface in the computer model does not measure synaptic change in a living organism.\par\nopagebreak[4]

\end{samepage}


\noindent\textit{Current registry lookup labels:} Biological plasticity; synaptic plasticity; biological learning\par

% Stable term ID: learning-before-the-final-task
\begin{samepage}
\subsection*{Learning before the final task}
\noindent Training first on another objective, such as predicting a later neural count, before training for a paired movement target.\par
\end{samepage}
\begin{samepage}
\noindent The paper calls this pretraining. Earlier source-task labels remain part of the information the model received, even when only new-setting labels are counted in a small budget.\par\nopagebreak[4]

\end{samepage}


\noindent\textit{Current registry lookup labels:} Pretraining; neural next-bin pretraining\par

% Stable term ID: mapping-movement-onto-another-body-model
\begin{samepage}
\subsection*{Mapping movement onto another body model}
\noindent Turn a specified movement quantity, such as fingertip velocity components, into movement goals for a particular computer body. The mapping depends on that body and the chosen task.\par
\end{samepage}
\noindent Mapping specified movement quantities into goals for a particular body model, with the source quantities, target coordinates and transformation declared. This does not by itself reconstruct full posture or demonstrate biological transfer.\par
\begin{samepage}
\noindent The retained source-specific adapter takes two macaque fingertip velocity components as its source and produces human-model shoulder and elbow goals. Their signs carry direction; scalar speed would discard it. These are goals for that model and task, so the mapping establishes neither cross-species skill transfer nor verified biological function.\par\nopagebreak[4]

\end{samepage}


\noindent\textit{Current registry lookup labels:} Retargeting; body adapter; organism-specific body adapter; joint goal\par

% Stable term ID: meme-ification
\begin{samepage}
\subsection*{Meme-jacking}
\noindent Borrow the recognition; earn the explanation. A familiar image should help the audience see the scientific point before the script develops it.\par
\end{samepage}
\noindent Use a recognizable meme or cultural reference to make a scientific point immediately legible, drawing on familiarity already earned by others.\par
\begin{samepage}
\noindent Memification hijacking is the user’s expansion. References may come from any online culture, not only films. Apply one purposeful reference per teaching concept without crowding formal methods or labels. Recognition does not grant reuse rights, endorsement, proven comprehension or scientific validity. Preserve the original Meme-ification history.\par\nopagebreak[4]

\end{samepage}
\noindent Meme-ification/memeification are spelling forms; mimification is retained only as the corrected historical transcription.\par


\noindent\textit{Current registry lookup labels:} Meme-ification; meme-ification; memeification; mimification; meme-jacking; meme jacking; memification hijacking\par

% Stable term ID: mind-meld
\subsection*{Mind-Meld}
\noindent What information could one system give another, and would it help with the task? Mind-Meld names the proposed first Brain2Brain encounter, whose signal and effects would have to be measured.\par
\noindent A proposed first Brain2Brain study of a specified information link between neural systems during a declared task.\par
\noindent The roadmap places this at step 3.1 before later NeuroTransfer studies. The name does not describe fused minds, shared private experience, arbitrary thought reading or demonstrated skill transfer. Specify the source, recipient, recording and feedback interfaces, task, controls and independent outcome. Assistance during a link and learning retained without trainer information need separate tests, with any necessary compensatory bridge held constant. A reported natural connection is a motivating case, not evidence that this engineered link exists.\par

\noindent\textit{Current registry lookup labels:} No additional label declared.\par

% Stable term ID: mind-merge
\subsection*{Mind-Merge}
\noindent An agent could guide practice by giving feedback about a stated task. Mind-Merge names the proposed study of that interaction: what help does the person receive, and what can they do later without it?\par
\noindent A proposed Agent2Brain study of task-specific feedback from a software agent to a person and its measurable effects.\par
\noindent The roadmap places this at step 4.1. Ordinary sensory feedback and direct neural delivery remain different routes with different interfaces, evidence and admission. EEG and EMG record; neither supplies a neural writer. No merged identity, shared mind, retained skill or physical intervention follows from the name. Preserve the reverse observation path, matched practice and delayed/new-task retention tests, removing trainer information while keeping a necessary bridge constant. Physical spinal work remains read-only.\par

\noindent\textit{Current registry lookup labels:} No additional label declared.\par

% Stable term ID: mixed-source-animal-and-added-tissue
\begin{samepage}
\subsection*{Mixed-source animal and added tissue}
\noindent A chimera has biological contributions from different sources. A graft is added tissue; when it crosses species, it is a xenotransplant. An organoid is an organized tissue model.\par
\end{samepage}
\begin{samepage}
\noindent Survival, physical connection, activity and useful function are separate questions. The proposed comparisons do not establish a humanized whole brain or an implemented treatment.\par\nopagebreak[4]

\end{samepage}


\noindent\textit{Current registry lookup labels:} Chimera, graft and organoid; chimera; chimerism; graft; organoid; host–graft; xenotransplant; xenograft; XenoTransplantation\par

% Stable term ID: motion-capture-composite-relay
\begin{samepage}
\subsection*{MoCap ReComp}
\noindent MoCap ReComp is Motion-Capture Relay-Composite: the proposed measurement component that brings views of the same movement into a common physical frame. It should preserve which view supplied each point, so the joined map can be compared with its original observations.\par
\noindent The proposal starts with two dot maps of the same moving body: one from above, one from below. Align their time and position in a common physical coordinate frame, then join them into one motion map while retaining the source of each point.\par
\end{samepage}
\noindent A proposed method that creates two dot maps (3D point clouds) of the same moving body, one captured from above and one from below, aligns them in time and a common physical coordinate frame, then stitches them into one composite motion map while retaining each point’s source.\par
\noindent To compare the same movement from above and below, the proposal must align both streams in time and a common physical coordinate frame while retaining each observation’s source. Bidirectional denotes those opposing measurements, not neural communication or causal direction. LiDAR ranging and infrared optical motion capture are proposed alternatives; marker triangulation and LiDAR point clouds are different measurements, so the actual modality and camera/ranging geometry must be qualified. The maze geometry calls for measurements perpendicular to the floor, with viewing paths from both sides through material intended to be clear. Transmission and material effects need testing; concentrated and distributed arrays remain alternatives. Hardware, coverage, timing, calibration, spatial alignment, occlusion, motion error and latency require validation against real references. The joined map is a derived record, not a direct observation of hidden joints, tissue mechanics, muscle commands or neural function. Neither successful implementation, complete body reconstruction nor material transparency has been established by the name. The single vertical camera illustration remains schematic.\par
\noindent\textit{Geometry:} Declare the top/bottom sensor positions, measured object, coordinate frame, viewing/illumination paths and coverage. Preserve the requested perpendicular-to-maze-floor geometry; document deviations rather than inferring them from an illustration.\par
\noindent\textit{Modality:} Declare LiDAR ranging or infrared optical motion capture, including marker-based/markerless method, camera count, triangulation geometry, material path and actual calibration. No complete reconstruction or physiological writing follows from opposing views.\par
\noindent\textit{Timing:} Record source timestamps, clocks, acquisition cadence, synchronization error and any interpolation, offline processing or end-to-end latency. No real-time capability is assumed.\par
\noindent\textit{Calibration and registration:} Preserve calibration data and transformations placing both streams into a common physical frame; measure registration error and keep original observations and reconstruction steps traceable.\par
\noindent\textit{Limits:} Record overlap, occlusion, missing observations, material effects, uncertainty and failure cases. Opposing measurements alone do not prove complete 3D pose, forces or biological function.\par
\noindent\textit{Physical scale:} Declare measured distance ranges and units for an actual acquisition; no supported numeric range or physical unit assignment is presently available for this proposal.\par
\noindent\textit{Statistical aggregation:} Declare any chosen evaluation grouping/averaging rule separately from physical scale and computational granularity.\par
\noindent\textit{References and acceptance:} Use \hyperref[app:reference-comparison-protocol]{the Reference Comparison Protocol and unfilled record} to identify the real-reference series needed for each acquisition component and the assembled map. Its provenance, rights, organism-stage, physical-unit, modality, clock, calibration, occlusion, contact, independent held-out, error, uncertainty and acceptance fields remain unfilled and HOLD. It supplies no measured series or apparatus validation. The unchanged programme reference criteria and every stronger original requirement govern their respective claims; no acceptance threshold is introduced or relaxed.\par
\noindent\textit{Outputs:} Retain the top dot map, bottom dot map and joined map as separate traceable outputs. Their timestamps and common-frame transformations must be explicit; do not fill unobserved regions as though they were measured.\par
\noindent\textit{Scope:} proposed application: Complementary top/bottom motion measurements of a body moving through the maze; concentrated coordinated platforms and distributed arrays remain alternatives; measurement boundary: Observed distances or geometry remain distinct from registered/stitched composites and any model-estimated pose or musculoskeletal quantity.\par
\noindent\textit{Organizational level:} whole system; signal/feature\par
\begin{samepage}
\noindent\textit{Computational granularity:} top dot map (3D point cloud); bottom dot map (3D point cloud); time alignment and spatial registration; joined composite motion sequence\par\nopagebreak[4]

\end{samepage}
\noindent Historical spellings name this same proposed capture method. Relay here means passing captured measurements onward; the separate agent-coordination Relay protocol does not describe motion capture.\par


\noindent\textit{Current registry lookup labels:} motion capture composite relay; motion-capture composite relay; motion capture composite; bidirectional composite motion capturing using LiDAR; Motion capture composite relay; motion capture relay composite; motion-capture relay composite; bidirectional LiDAR; opposed-view infrared motion capture; bidirectional infrared; motion-capture relay-composite; Motion-Capture Relay-Composite\par

% Stable term ID: model-that-chooses-a-label
\begin{samepage}
\subsection*{Model that chooses a label}
\noindent A statistical model that assigns scores to categories, such as the ten attempted-finger labels.\par
\end{samepage}
\begin{samepage}
\noindent The selected human model uses logistic classification, but its scores do not establish validated confidence in biological truth. Its reported accuracy comes from development data that have already been seen.\par\nopagebreak[4]

\end{samepage}


\noindent\textit{Current registry lookup labels:} Logistic classifier; logistic regression; classifier; classification; logistic\par

% Stable term ID: model-weights-kept-fixed
\begin{samepage}
\subsection*{Model weights kept fixed}
\noindent Keeping the fitted numbers in a model unchanged during a defined comparison.\par
\end{samepage}
\begin{samepage}
\noindent The tested transfer trains a connector while retaining the older main model. Fixed weights make the comparison clear; they do not prove biological sameness across days.\par\nopagebreak[4]

\end{samepage}


\noindent\textit{Current registry lookup labels:} Frozen parameters; frozen trunk; frozen readout; frozen checkpoint\par

% Stable term ID: modeled-route-from-one-muscle-attachment-to-another
\begin{samepage}
\subsection*{Modeled route from one muscle attachment to another}
\noindent The route and length assigned to a muscle and tendon as the model’s joints move.\par
\end{samepage}
\begin{samepage}
\noindent The native arm retains 39 named routes. Their modeled lengths differ from direct measurements of living tissue, muscle commands and force.\par\nopagebreak[4]

\end{samepage}


\noindent\textit{Current registry lookup labels:} Muscle-tendon path; musculotendon length; modeled muscle length; geometric path\par

% Stable term ID: models-that-convert-inputs-or-predict-outputs
\begin{samepage}
\subsection*{Models that convert inputs or predict outputs}
\noindent One model can turn an input into a representation; another can use available inputs to predict a named result.\par
\end{samepage}
\begin{samepage}
\noindent The paper calls these encoders and decoders. Predicting counts from mechanics or mechanics from counts does not by itself show the biological cause of either.\par\nopagebreak[4]

\end{samepage}


\noindent\textit{Current registry lookup labels:} Encoder and decoder; encoding; decoding; readout; decode; encoder; decoder\par

% Stable term ID: movement-brain-area
\begin{samepage}
\subsection*{Movement brain area}
\noindent A brain area involved in movement. The paper uses the names motor cortex, MC and M1 for relevant source recordings.\par
\end{samepage}
\begin{samepage}
\noindent The signals help predict recorded movement-related targets. Knowing the brain area does not identify matching cells across recordings or reveal a whole-body command pathway.\par\nopagebreak[4]

\end{samepage}


\noindent\textit{Current registry lookup labels:} Motor cortex; M1; MC; primary motor cortex\par

% Stable term ID: muscle-electrical-recording
\begin{samepage}
\subsection*{Muscle electrical recording}
\noindent Electrical activity recorded from muscles; the method is called EMG.\par
\end{samepage}
\begin{samepage}
\noindent Muscle electrical activity differs from force and from the computer model’s command and activity states. The trained targets here are rescaled signal summaries.\par\nopagebreak[4]

\end{samepage}


\noindent\textit{Current registry lookup labels:} Electromyography; EMG\par

% Stable term ID: myoception
\subsection*{Myoception}
\noindent Sensing muscle-driven movement and the feedback that comes with it.\par
\noindent The programme’s term for the sensory side of musculoskeletal contraction and movement, including proprioception, tactile feedback and other task-relevant sensations linked to action.\par
\noindent This is the user’s proposed programme scope, not a claim that every relevant signal has been recorded or a new external scientific standard established. Muscle electrical activity, proprioceptive or tactile measurements, neural recordings and simulated feedback remain distinct observations. EMG does not directly record sensory afferents or the full sensory consequences of contraction. Declare each source’s organism, anatomy, modality, task, units, timing, uncertainty and measured, derived or simulated status. The wider neural-data ambition includes brain, spinal and peripheral activity and relevant visceral signals; the current focus is the motor–sensory system. No body-wide coverage, unique hidden neural state or subjective experience is established by this term.\par

\noindent\textit{Current registry lookup labels:} myoception; M-Y-O-ception\par

% Stable term ID: musculoskeletal-simulation-physics-engine
\begin{samepage}
\subsection*{MyoSim Replay}
\noindent MyoSim Replay is the proposed MusculoSkeletal Simulation Physics-Engine. Declared anatomy, mechanical parameters and control inputs would produce movement predictions to compare with independent measurements. The external MyoSim body-model collection is a separate resource. Simulation replay is not a measured biological replay.\par
\noindent The proposed digital model would calculate how muscles and the skeleton move and exert forces, with real measurements used to check its predictions.\par
\end{samepage}
\noindent The proposed physics-based simulation component that models specified muscles, bones, joints and their mechanical interactions, using declared anatomy, parameters, forces and contact rules, with predictions tested against independent movement and other applicable physical measurements.\par
\noindent In the proposed five-component arrangement, declared anatomy, muscle/control inputs and physical context would drive a simulation of movement and forces. The motion-capture relay-composite would provide traceable measured movement to compare with those predictions. That comparison needs matching timing, coordinates, units and calibration, and tests that preserve anatomical and physical constraints, uncertainty, held-out cases, baseline controls, a series of real references and independent acceptance. Movement geometry alone does not measure muscle activation, internal forces, hidden neural state or subjective experience; force, activation and causal claims need additional independent references. Visual effects (VFX) rendering shows the model but cannot validate its mechanics. This remains a proposed simulation component, not a newly demonstrated implementation or external scientific field. Approximate simulation and biological virtualization verified within declared error bounds remain distinct; this component establishes neither universal precision nor whole-body completeness.\par
\noindent\textit{Scope:} programme role: Virtual-body mechanics component; initially constrained hand and foot tasks; inputs and outputs: Declared anatomical parameters, muscle/control inputs and physical context produce simulated movement and mechanical quantities. Independently acquired motion and other necessary observables remain references, not simulator ground truth.\par
\noindent\textit{Organizational level:} whole system; organ\par
\begin{samepage}
\noindent\textit{Computational granularity:} musculoskeletal model; mechanical solver; measurement-to-prediction comparison\par\nopagebreak[4]

\end{samepage}


{\rightskip=0pt plus1fil\relax
\noindent\textit{Current registry lookup labels:} musculoskeletal simulation physics engine; musculoskeletal-simulation physics engine; musculoskeletal simulation; musculoskeletal physics model; musculoskeletal-simulation physics-engine; MusculoSkeletal Simulation Physics-Engine\par
}

% Stable term ID: names-used-for-source-subjects
\begin{samepage}
\subsection*{Names used for source subjects}
\noindent Indy, Han and Monkey L label macaque sources; N1 labels the human finger source.\par
\end{samepage}
\begin{samepage}
\noindent These labels keep different subjects separate. Similar brain-region or task names do not make them one synchronized recording.\par\nopagebreak[4]

\end{samepage}


\noindent\textit{Current registry lookup labels:} Source subject labels; Indy; Han; Monkey L; N1\par

% Stable term ID: nerves-bringing-body-information-in
\begin{samepage}
\subsection*{Nerves bringing body information in}
\noindent Afferent paths carry body-related information toward the brain and spinal cord, together called the central nervous system.\par
\end{samepage}
\begin{samepage}
\noindent For example, a limb-related sensory signal can travel inward. The source-specific cortical predictions do not measure this whole path or make missing nerve activity observed.\par\nopagebreak[4]

\end{samepage}


\noindent\textit{Current registry lookup labels:} afferent; afferent pathway; body-to-CNS; CNS; central nervous system\par

% Stable term ID: nerves-carrying-commands-out
\begin{samepage}
\subsection*{Nerves carrying commands out}
\noindent Efferent paths carry outputs from the brain and spinal cord toward the body, including muscles and organs.\par
\end{samepage}
\begin{samepage}
\noindent A prediction of movement is not a recording of that outward pathway. The prototype does not supply a synchronized voluntary sequence through spinal cord, nerves and muscles.\par\nopagebreak[4]

\end{samepage}


\noindent\textit{Current registry lookup labels:} efferent; efferent pathway; CNS-to-body\par

% Stable term ID: neural-decoding
\begin{samepage}
\subsection*{Neural decoding}
\noindent Estimating declared information such as intended action from recorded neural observations using a specified model.\par
\end{samepage}
\noindent Decode accuracy, channels, source identity, task and held-out domain must be demonstrated. Brain-only control requires exclusion of covert muscle, eye and motion signals; decoding does not write signals into the brain.\par
\noindent\textit{Organizational level:} Declared game/model/interface role\par
\noindent\textit{Computational granularity:} Declared task-specific model or interface\par
\noindent\textit{Statistical aggregation:} Validation scope declared per experiment; no physical-scale inference\par
\noindent\textit{Source:} Restoring cortical control of functional movement in a human with quadriplegia\par
\noindent\textit{Source identifier:} bouton2016\par
\noindent\textit{DOI:} 10.1038/nature17435\par
\noindent\url{https://www.nature.com/articles/nature17435}\par
\noindent Research precedent linking implanted cortical recording to a custom neuromuscular stimulation sleeve; not an EEG/OpenBCI or elite-skill result\par
\noindent\textit{Source:} A brain-computer interface that evokes tactile sensations improves robotic arm control\par
\noindent\textit{Source identifier:} flesher2021\par
\noindent\textit{DOI:} 10.1126/science.abd0380\par
\noindent\url{https://pubmed.ncbi.nlm.nih.gov/34016775/}\par
\noindent Research precedent for motor decoding plus sensory feedback through intracortical microstimulation; invasive and not a consumer EEG capability\par


\noindent\textit{Current registry lookup labels:} No additional label declared.\par

% Stable term ID: neural-encoding
\begin{samepage}
\subsection*{Neural encoding}
\noindent Mapping declared information to a representation in neural activity; in a sensory-interface design, constructing candidate stimulation patterns for specified feedback.\par
\end{samepage}
\noindent A computational encoder is distinct from a qualified stimulation device and evidence that a participant receives the intended sensation. The term does not mean arbitrary thought writing or skill upload.\par
\noindent\textit{Organizational level:} Declared game/model/interface role\par
\noindent\textit{Computational granularity:} Declared task-specific model or interface\par
\noindent\textit{Statistical aggregation:} Validation scope declared per experiment; no physical-scale inference\par
\noindent\textit{Source:} A brain-computer interface that evokes tactile sensations improves robotic arm control\par
\noindent\textit{Source identifier:} flesher2021\par
\noindent\textit{DOI:} 10.1126/science.abd0380\par
\noindent\url{https://pubmed.ncbi.nlm.nih.gov/34016775/}\par
\noindent Research precedent for motor decoding plus sensory feedback through intracortical microstimulation; invasive and not a consumer EEG capability\par


\noindent\textit{Current registry lookup labels:} No additional label declared.\par

% Stable term ID: playing-back-a-recording
\begin{samepage}
\subsection*{Neural Replay}
\noindent Revisiting a stored sequence of neural samples with its stated timing and channel order.\par
\end{samepage}
\begin{samepage}
\noindent This playback is not a live recording or instructions for stimulating a nervous system. A displayed arm movement does not show what the source animal’s muscles actually did.\par\nopagebreak[4]

\end{samepage}


\noindent\textit{Current registry lookup labels:} Neural replay; recorded replay; archived replay; cached playback; Playing back a recording\par

% Stable term ID: neuroagent-tutor
\subsection*{NeuroAgent Trainer}
\noindent A software agent would guide a trainee during a stated task.\par
\noindent A proposed software-agent trainer role in a declared Agent2Brain NeuroTransfer study.\par
\noindent Identify the actual software model and feedback route. Computed activations/events do not make a biological trainee digital. Illustrative connections do not establish direct neural writing, successful transfer or retained learning. NeuroAgent Tutor is a historical label; NeuroAgent Trainer is the current name. It describes a proposed software-agent trainer role, without selecting the agent or qualifying its feedback interface.\par
\noindent\textit{Organizational level:} whole system; signal/feature\par

\noindent\textit{Current registry lookup labels:} NeuroAgent Tutor\par

% Stable term ID: neuroarcher
\begin{samepage}
\subsection*{NeuroArcher}
\noindent A proposed archery game that studies coordinated use of both hands and arms through physical motion control, brain-signal control of a digital body, and a separately admitted later assisted-body experiment.\par
\end{samepage}
\noindent The near-term prop is a toy bow with small suction-cup tips and an over-arm arrow pouch. Target/creature choice is unresolved. Simulated performance does not establish muscle stimulation or retained biological archery skill.\par
\noindent\textit{Organizational level:} Declared game/model/interface role\par
\noindent\textit{Computational granularity:} Declared task-specific model or interface\par
\begin{samepage}
\noindent\textit{Statistical aggregation:} Validation scope declared per experiment; no physical-scale inference\par\nopagebreak[4]

\end{samepage}


\noindent\textit{Current registry lookup labels:} Neuro Archer\par

% Stable term ID: neuroboxer
\begin{samepage}
\subsection*{NeuroBoxer}
\noindent A proposed boxing game that compares arm-motion control, brain-signal control of a digital body, and a separately admitted later assisted-body experiment.\par
\end{samepage}
\noindent Prime-era Mike Tyson is the requested virtual opponent and meme reference. A game, recorded prediction, stimulation-assisted motion and retained biological skill learning need different evidence. No actual opponent, head contact, finished game or stimulation capability is claimed.\par
\noindent\textit{Organizational level:} Declared game/model/interface role\par
\noindent\textit{Computational granularity:} Declared task-specific model or interface\par
\begin{samepage}
\noindent\textit{Statistical aggregation:} Validation scope declared per experiment; no physical-scale inference\par\nopagebreak[4]

\end{samepage}


\noindent\textit{Current registry lookup labels:} Neuro Boxer\par

% Stable term ID: neuroconnectome
\begin{samepage}
\subsection*{NeuroConnectome}
\noindent A map of specified neural connections used to constrain a model of the nervous system.\par
\end{samepage}
\begin{samepage}
\noindent Connectome is the established term. The user’s NeuroConnectome label does not establish complete coverage or identify functional dynamics, state, plasticity, sensory context or unique thought from wiring alone. Declare species, anatomical coverage, resolution and provenance.\par\nopagebreak[4]

\end{samepage}


\noindent\textit{Current registry lookup labels:} Neuroconnectome\par

% Stable term ID: neurohandler
% Canonical status: "user_requested_proposed_programme_role"
\subsection*{NeuroHandler}
\noindent How well can a hand carry out the chosen task, and what information helps it do so? NeuroHandler is the proposed manual-dexterity experiment, with the hand’s measured action providing the result to explain.\par
\noindent A proposed manual-dexterity experiment studying specified neural/body observations, task control and independently measured outcomes.\par
\noindent The earlier proposed manual tasks remain: rapid hand rubbing, the requested knuckle gesture and the interlaced-hand stretch. NeuroHandler names the manual-dexterity experiment; it does not establish a finished apparatus or successful model. Define the action, available inputs, physical or simulated setting, comparison and outcome before evaluation. Keep motion control, neural prediction, assistance and retained learning separate; any physical output requires separate qualification and admission. No dexterity gain, clinical benefit or intervention procedure is established.\par


\noindent\textit{Current registry lookup labels:} No additional label declared.\par

% Stable term ID: neurojuggler
% Canonical status: "user_requested_proposed_programme_role"
\subsection*{NeuroJuggler}
\noindent Moving each limb successfully does not settle whether the limbs work together. NeuroJuggler is the proposed coordination experiment: define the joint task, then measure the relationship between the limb actions and the result.\par
\noindent A proposed multi-limb coordination experiment studying how specified limb actions are coordinated in a declared task.\par
\noindent The earlier coordination proposal remains a hot-potato transfer between the hands, followed preferably by two-ball juggling. Those are proposed tasks; the final apparatus and qualified measurement setup remain unestablished. The label supplies no successful performance or implemented neural interface. Declare the task, limbs, inputs, timing, body model or physical setting and independent outcomes. Separate coordination from isolated-limb performance, simulated success from biological function, and current assistance from retained learning. Physical output, safety, rights and held-out tests keep their own admission requirements.\par


\noindent\textit{Current registry lookup labels:} No additional label declared.\par

% Stable term ID: neuromuscular-electrical-stimulation
\begin{samepage}
\subsection*{Neuromuscular electrical stimulation}
\noindent Electrical stimulation intended to elicit muscle contraction through the relevant neuromuscular pathways.\par
\end{samepage}
\noindent NMES is a writing/actuation intervention, whereas EMG records muscle activity. EMG electrodes or an OpenBCI recording board do not by themselves provide qualified stimulation. Neither contraction nor imposed movement demonstrates retained skill.\par
\noindent\textit{Organizational level:} Declared game/model/interface role\par
\noindent\textit{Computational granularity:} Declared task-specific model or interface\par
\begin{samepage}
\noindent\textit{Statistical aggregation:} Validation scope declared per experiment; no physical-scale inference\par\nopagebreak[4]

\end{samepage}
\noindent\textit{Source:} Restoring cortical control of functional movement in a human with quadriplegia\par
\noindent\textit{Source identifier:} bouton2016\par
\noindent\textit{DOI:} 10.1038/nature17435\par
\noindent\url{https://www.nature.com/articles/nature17435}\par
\noindent Research precedent linking implanted cortical recording to a custom neuromuscular stimulation sleeve; not an EEG/OpenBCI or elite-skill result\par


\noindent\textit{Current registry lookup labels:} NMES; Electrical muscle stimulation; EMS\par

% Stable term ID: neuroreplay
\subsection*{NeuroReplay}
\noindent What changes when a trainee receives replay-based guidance? NeuroReplay names that proposed study: compare performance while guidance is available, then test what remains when the guidance ends.\par
\noindent A proposed Stage 0 research direction for studying replay-guided sensorimotor feedback and its effects on current performance and later learning.\par
\noindent This is the new roadmap label for steps 0.2–0.5, with a derived working scope; it does not rename Neural Replay, the existing stored-sample playback entry playing-back-a-recording. A stored sequence can supply a comparison, but playback alone establishes neither delivery to a nervous system nor biological learning. Fix the task, signal, feedback route, timing and equal-practice comparisons; distinguish matched, stale or unmatched replay, no-link and inactive or sham controls. Test delayed retention and new tasks after removing trainer information while keeping any necessary compensatory bridge constant. The questioned antechinus age ranges are not established developmental or intervention windows. No physical writer, procedure or present result is supplied by this label.\par

\noindent\textit{Current registry lookup labels:} No additional label declared.\par

% Stable term ID: neuroscience-recording-files
\begin{samepage}
\subsection*{Neuroscience recording files}
\noindent NWB stores neuroscience data and descriptions. In the paired source files here, SET holds recording descriptions and FDT holds binary samples.\par
\end{samepage}
\begin{samepage}
\noindent A stale unit in the NWB muscle source must not override verified processing. SET headers match source hashes; FDT previews have checked local byte ranges rather than full verification of each upstream file.\par\nopagebreak[4]

\end{samepage}


\noindent\textit{Current registry lookup labels:} Neurodata Without Borders file; NWB; Recording header and sample file; SET and FDT files; SET; FDT; EEGLAB header; byte-range preview\par

% Stable term ID: neurotransfer
\begin{samepage}
\subsection*{NeuroTransfer}
\noindent Specify what would move between which systems, then test what the recipient can do with and without assistance.\par
\end{samepage}
\noindent The proposed transfer of specified task-relevant information or a mapping between declared systems, with assistance and retained learning evaluated separately.\par
\begin{samepage}
\noindent The label does not establish a successful biological transfer or retained skill. A weight copy, recorded replay, guided movement and independently retained learning answer different questions. Each proposed interface needs its own source, target, task, controls, rights and qualified outcome.\par\nopagebreak[4]

\end{samepage}


\noindent\textit{Current registry lookup labels:} neurotransfer\par

% Stable term ID: neurotranslate
\begin{samepage}
\subsection*{NeuroTranslate}
\noindent Brain2Body Neural Activity is the long-form component label paired with NeuroTranslate in this proposed architecture. Recorded activity is an observation. NeuroTranslate names the mapping operation that could relate specified observations to movement, or reconstruct specified recorded features in the reverse direction. The observation, learned relationship and predicted value remain different things.\par
\noindent State the observations, the target and the task, then test how well the mapping works.\par
\end{samepage}
\noindent The proposed operation of mapping specified neural and body observations between declared inputs and targets in a stated task.\par
\begin{samepage}
\noindent A model could predict movement from specified observations, or reconstruct specified recorded features in the reverse direction. Each direction needs its own test in the stated task. Neither establishes a unique hidden state, sensory delivery or biological skill transfer. NeuroTranslate names this proposed mapping operation; NeuroTranslation names the field and NeuroTranslator the implemented software.\par\nopagebreak[4]

\end{samepage}


\noindent\textit{Current registry lookup labels:} neurotranslate; Brain2Body Neural Activity\par

% Stable term ID: neurotranslation
\begin{samepage}
\subsection*{NeuroTranslation}
\noindent How could a particular neural recording be related to a defined movement task? NeuroTranslation is the proposed study of mapping nerve signals and the patterns that describe them between clearly defined situations like these. Every mapping must state where its inputs come from, what its outputs mean, and how it is tested.\par
\end{samepage}
\noindent A proposed field for translating or mapping neural signals and representations between explicitly defined source and target domains, with stated context, scale, tests and limitations.\par
\noindent Each connection must be assessed against its named source and stated test. The project has not demonstrated a universal brain language or an uploaded skill.\par
\noindent\textit{Scope:} macro: Relationships between neural systems, body systems and defined tasks or organisms; micro: Mappings involving identified local signals, channels, features and circuit-level representations.\par
\begin{samepage}
\noindent\textit{Organizational level:} whole system; network; organ; local circuit; signal/feature\par\nopagebreak[4]
\nopagebreak[4]
\end{samepage}


\noindent\textit{Current registry lookup labels:} neurotranslation; neural translation; Connecting neural recordings to defined tasks; Neurotranslation\par

% Stable term ID: source-conditioned-neurotranslation-architecture
\begin{samepage}
\subsection*{NeuroTranslation Network}
\noindent A plan for connecting more neural and body components while keeping the origin and meaning of each value clear.\par
\end{samepage}
\begin{samepage}
\noindent The proposed network would connect neural and body components while carrying the origin and meaning of each value through the mapping. Each link could then be examined in its own recording context. Body-wide translation and biological teaching remain unproved.\par\nopagebreak[4]

\end{samepage}


\noindent\textit{Current registry lookup labels:} Proposed larger design; Source-Conditioned Neurotranslation Architecture; SCNA; Neurotranslation Network\par

% Stable term ID: neurotranslation-engine
\begin{samepage}
\subsection*{NeuroTranslator}
\noindent The software that joins particular recordings, prediction models, playback and computer-body checks.\par
\end{samepage}
\begin{samepage}
\noindent Predictions and playback can be checked against the particular recordings they use. That is the scope of the implemented NeuroTranslator; neither a universal brain language nor skill upload has been demonstrated.\par\nopagebreak[4]

\end{samepage}


\noindent\textit{Current registry lookup labels:} Working translation software; Neurotranslation Engine; NTE; Neurotranslator\par

% Stable term ID: neurotwin-tutor
\subsection*{NeuroTwin Trainer}
\noindent A biological trainer would guide a trainee during a stated task.\par
\noindent A proposed biological trainer role in a declared Brain2Brain NeuroTransfer study.\par
\noindent The requested name alone establishes neither genetic twins, a validated digital twin, communication nor retained learning. Biological electrical activity is electrochemical and can include graded voltages and spikes; identify any actual measurement, otherwise keep the activity illustrative. NeuroTwin Tutor is a historical label; NeuroTwin Trainer is the current name. These names belong to a proposed programme view. They establish neither an implemented interface nor a demonstrated learning result.\par
\noindent\textit{Organizational level:} whole system; signal/feature\par

\noindent\textit{Current registry lookup labels:} NeuroTwin Tutor\par

% Stable term ID: numerical-arrays-saved-arrays-and-prediction
\begin{samepage}
\subsection*{Numerical arrays, saved arrays and prediction}
\noindent NumPy provides tools for numerical arrays. NPZ saves arrays; inference means running a trained model to produce an output.\par
\end{samepage}
\begin{samepage}
\noindent Portable execution still needs the correct processing and source identity. Matching saved-model outputs does not establish success on new organisms or biological skills.\par\nopagebreak[4]

\end{samepage}


\noindent\textit{Current registry lookup labels:} NumPy inference and NPZ archive; NumPy; NPZ; portable inference; inference\par

% Stable term ID: numerical-checks-on-allowed-movement
\begin{samepage}
\subsection*{Numerical checks on allowed movement}
\noindent Checks that the computer model stays defined and within its stated joint limits.\par
\end{samepage}
\begin{samepage}
\noindent Joint-limit penetration means going past an angular bound, not entering tissue. The 0.001 rad tolerance is a numerical rule, not a certificate of correct neural commands.\par\nopagebreak[4]

\end{samepage}


\noindent\textit{Current registry lookup labels:} Mechanical feasibility and anatomical tolerance; anatomical limit; limit penetration; finite state; joint range; numerical tolerance\par

% Stable term ID: numerical-software-and-its-data-file
\begin{samepage}
\subsection*{Numerical software and its data file}
\noindent MATLAB is numerical-computing software, and MAT is a format for its structured data.\par
\end{samepage}
\begin{samepage}
\noindent The second Indy recording uses the original MAT source. The format does not prove matching cells, past-only source processing or complete observation.\par\nopagebreak[4]

\end{samepage}


\noindent\textit{Current registry lookup labels:} MATLAB and MAT file; MATLAB; MAT\par

% Stable term ID: observe-it-connect-it-keep-it-useful
\begin{samepage}
\begin{samepage}
\subsection*{Observe it, connect it, keep it useful}
\noindent Observation burden, prediction on excluded cases and lasting task benefit need separate measurements. A clear picture does not answer those experimental questions.\par\nopagebreak[4]

\end{samepage}
\end{samepage}


\noindent\textit{Current registry lookup labels:} Access it, Ground it, Keep it; access; grounding; stability; signal stability\par

% Stable term ID: one-multiplier-for-modeled-lengths
\begin{samepage}
\subsection*{One multiplier for modeled lengths}
\noindent A single number chosen to reduce differences between the model’s muscle lengths and the source-derived length references.\par
\end{samepage}
\begin{samepage}
\noindent The number is fitted on training poses and checked on distinct observed poses. It is not complete body calibration; no movement-control experiment is run on the scaled candidate.\par\nopagebreak[4]

\end{samepage}


\noindent\textit{Current registry lookup labels:} Scalar length fit; global length factor; uniform factor; a*; scalar bridge; recording length-fit candidate; SDK-scaled candidate\par

% Stable term ID: one-recording-period
\begin{samepage}
\subsection*{One recording period}
\noindent One experiment or recording period with its own subject, channels, timing and task.\par
\end{samepage}
\begin{samepage}
\noindent Two recording days, or recordings from different animals, do not form one observed nervous system simply because the files are loaded together.\par\nopagebreak[4]

\end{samepage}


\noindent\textit{Current registry lookup labels:} Recording and session; source recording; session; recording family\par

% Stable term ID: open-body-model-collection
\begin{samepage}
\subsection*{Open body-model collection}
\noindent MyoSim supplies computer-body models with muscles and bones, including models used in this project.\par
\end{samepage}
\begin{samepage}
\noindent The collection is earlier work, not a new project invention. Availability of an arm model does not establish matching macaque and human muscles or neural commands.\par\nopagebreak[4]

\end{samepage}


\noindent\textit{Current registry lookup labels:} MyoSim, MyoSuite and MyoArm; MyoSim\par

% Stable term ID: open-muscle-control-task-tools
\begin{samepage}
\subsection*{Open muscle-control task tools}
\noindent MyoSuite provides tasks and tools for experiments with muscle-driven computer-body control.\par
\end{samepage}
\begin{samepage}
\noindent It is earlier software reused here. Success on a simulated task is not proof of biological learning or a recovered neural command.\par\nopagebreak[4]

\end{samepage}


\noindent\textit{Current registry lookup labels:} MyoSuite\par

% Stable term ID: other-published-models-of-neural-recordings
\begin{samepage}
\subsection*{Other published models of neural recordings}
\noindent NDT2 learns patterns across recordings made inside the cortex, the brain’s outer layer. POYO handles differing sets of recorded neural features by representing events and learning relationships between them.\par
\end{samepage}
\begin{samepage}
\noindent The paper cites these as earlier work. It does not claim to have implemented them or use their existence as proof of universal transfer or matching cells.\par\nopagebreak[4]

\end{samepage}


\noindent\textit{Current registry lookup labels:} NDT2; large-scale neural pretraining; Published model for varying recorded populations; POYO; neural-event tokenization; cross-attention bottleneck\par

% Stable term ID: passing-information-and-gaining-a-lasting-skill
\begin{samepage}
\subsection*{Passing information and gaining a lasting skill}
\noindent Receiving useful task information differs from acquiring a durable ability that works independently in new conditions.\par
\end{samepage}
\begin{samepage}
\noindent A short software label budget is not measured biological practice time. Sewing, cube solving and breakdancing remain future examples without matched skill recordings here.\par\nopagebreak[4]

\end{samepage}


\noindent\textit{Current registry lookup labels:} Task information transfer and skill acquisition; task transfer; skill transfer; skill upload; compressed acquisition; learning speed; learning; physical skill learning; skill acquisition\par

% Stable term ID: past-only-inputs-and-processing-that-uses-later-data
\begin{samepage}
\subsection*{Past-only inputs and processing that uses later data}
\noindent The local model uses completed past inputs. Some earlier source processing may still have used later samples.\par
\end{samepage}
\begin{samepage}
\noindent Past-only indexing cannot reverse earlier filtering or rescaling that used the whole recording. It does not certify the full measurement-to-output delay in a living system.\par\nopagebreak[4]

\end{samepage}


\noindent\textit{Current registry lookup labels:} Causal indexing and acausal preprocessing; causal; algorithmic support; acausal filtering; offline interpolation; upstream filtering delay\par

% Stable term ID: peripheral-nerve-electrical-recording
\begin{samepage}
\subsection*{Peripheral nerve electrical recording}
\noindent A method for recording electrical signals associated with peripheral nerves, called ENG.\par
\end{samepage}
\begin{samepage}
\noindent The human preview contains aggregate electrical observations. It does not provide separately identified nerve-cell spike trains or a complete voluntary command sequence.\par\nopagebreak[4]

\end{samepage}


\noindent\textit{Current registry lookup labels:} Electroneurography; ENG; surface ENG; peripheral recording\par

% Stable term ID: phenotype
% Canonical status: "established_scientific_reference"
\subsection*{Phenotype}
\noindent The trait is what must be observed: for example, a specified body characteristic or response in the setting being studied. Phenotype names these observable traits, rather than the genetic description alone.\par
\noindent The observable traits of an organism arising through genotype, environment and development.\par
\noindent State the actual trait, measurement method, developmental and environmental context. A phenotype is not a universal function that can be read uniquely from structure or sequence. Retain the distinction between a measured trait and a model’s estimate.\par


\noindent\textit{Current registry lookup labels:} No additional label declared.\par

% Stable term ID: placing-body-parts-from-joint-positions
\begin{samepage}
\subsection*{Placing body parts from joint positions}
\noindent Calculating where body parts and muscle routes lie when joint positions are given.\par
\end{samepage}
\begin{samepage}
\noindent The formal name is forward kinematics, or FK. It places geometry; it does not calculate muscle-driven motion through time or recover muscle commands.\par\nopagebreak[4]

\end{samepage}


\noindent\textit{Current registry lookup labels:} Forward kinematics; FK; kinematic replay\par

% Stable term ID: planned-commands-corrections-and-constant-input
\begin{samepage}
\subsection*{Planned commands, corrections and constant input}
\noindent A planned command is feedforward; adjusting it from current model information is feedback; a constant muscle input is tonic control.\par
\end{samepage}
\begin{samepage}
\noindent The human correction can change the bounded command goal. Recorded movement stays a scoring reference, and the native correction candidates remain rejected against constant-input comparisons.\par\nopagebreak[4]

\end{samepage}


\noindent\textit{Current registry lookup labels:} Feedforward, feedback and tonic control; feedforward; feedback; corrected controller; tonic excitation; constant tonic\par

% Stable term ID: prediction-score-compared-with-a-mean-value-guess
\begin{samepage}
\subsection*{Prediction score compared with a mean-value guess}
\noindent R² compares squared prediction errors with a constant mean-value reference. Higher values usually indicate better prediction on the stated data.\par
\end{samepage}
\begin{samepage}
\noindent The score can be negative. Different recordings, target processing and ways of combining outputs prevent treating all R² values as directly comparable.\par\nopagebreak[4]

\end{samepage}


\noindent\textit{Current registry lookup labels:} Coefficient of determination; R²; R2; R\textasciicircum{}2; variance explained\par

% Stable term ID: processing-within-spinal-and-peripheral-pathways
\begin{samepage}
\subsection*{Processing within spinal and peripheral pathways}
\noindent Body-related information can be processed locally in spinal and peripheral structures rather than passing through one simple unchanged wire.\par
\end{samepage}
\begin{samepage}
\noindent A cortex-only prediction does not observe those local processes. The separate recordings and computer-body checks do not establish a complete measured brain-to-body causal chain.\par\nopagebreak[4]

\end{samepage}


\noindent\textit{Current registry lookup labels:} local spinal processing; local peripheral processing; spinal/peripheral processing\par

% Stable term ID: programmable-neurobiology
\begin{samepage}
\subsection*{Programmable neurobiology}
\noindent The proposed ability to specify a bounded nervous-system or sensorimotor outcome, select a qualified intervention, predict its effects, measure what happens and adapt within a validated closed loop.\par
\end{samepage}
\noindent\textit{Scope:} Brain, spinal, peripheral, muscle, mechanical and sensory feedback roles need separate evidence. Bidirectional prediction is different from safe sensory/action writing. An intermediary may improve a mapping without making it one-to-one. Arbitrary skill upload or universal disease/disability repair is not demonstrated.\par
\noindent\textit{Organizational level:} Declared system or interface\par
\noindent\textit{Computational granularity:} Declared process/interface/model\par
\noindent\textit{Statistical aggregation:} Validation domain stated independently of physical scale\par


\noindent\textit{Current registry lookup labels:} No additional label declared.\par

% Stable term ID: programmable-synthetic-biology
\begin{samepage}
\subsection*{Programmable synthetic biology}
\noindent Engineering specified biological behaviours by combining design, context-dependent molecular prediction, experimental measurement and iterative correction.\par
\end{samepage}
\noindent DNA/regulation, RNA/cell state and protein/interactions are explanatory layers, not a sufficient literal three-model recipe. Metabolism, feedback, environment and measured function matter. AlphaGenome and AlphaFold do not jointly certify whole-cell programmability; AlphaTranscriptome is a desired capability placeholder here.\par
\noindent\textit{Organizational level:} Declared system or interface\par
\noindent\textit{Computational granularity:} Declared process/interface/model\par
\begin{samepage}
\noindent\textit{Statistical aggregation:} Validation domain stated independently of physical scale\par\nopagebreak[4]

\end{samepage}


\noindent\textit{Current registry lookup labels:} Programmable molecular biology\par

% Stable term ID: public-reaching-recordings
\begin{samepage}
\subsection*{Public reaching recordings}
\noindent Public files pairing cortical electrical events with reaching behavior. The sources use names such as RTT and MC\_RTT.\par
\end{samepage}
\begin{samepage}
\noindent The original model’s release and the separately acquired Indy day have different processing. Their dates and target definitions must remain explicit.\par\nopagebreak[4]

\end{samepage}


\noindent\textit{Current registry lookup labels:} Reaching source archive; RTT; MC\_RTT; MC\textbackslash{}\_RTT; O'Doherty reaching archive; Sabes reaching corpus\par

% Stable term ID: public-source-records-and-document-links
\begin{samepage}
\subsection*{Public source records and document links}
\noindent Repositories such as DANDI and OpenNeuro identify source versions. PDF is a document format, and a URL gives a resource location.\par
\end{samepage}
\begin{samepage}
\noindent Source licenses and credit must be retained. A local paper draft or public input dataset does not mean the project’s new code is published, licensed or submitted to arXiv or a journal.\par\nopagebreak[4]

\end{samepage}


\noindent\textit{Current registry lookup labels:} Public records, document files and links; DANDI; OpenNeuro; PDF; URL; arXiv; public deposition; license; attribution\par

% Stable term ID: published-models-for-brain-scans-and-scalp-recordings
\begin{samepage}
\subsection*{Published models for brain scans and scalp recordings}
\noindent BrainLM learns from fMRI brain-scan data. Neuro-GPT combines an EEG input model with a model of changes over time and is discussed for motor imagery.\par
\end{samepage}
\begin{samepage}
\noindent These are earlier methods for different measurements. The project does not claim to have trained them, make scans and electrical spikes interchangeable, or upload skills.\par\nopagebreak[4]

\end{samepage}


\noindent\textit{Current registry lookup labels:} BrainLM; Published model of scalp signals; Neuro-GPT; NeuroGPT\par

% Stable term ID: published-neural-model-comparison-tasks
\begin{samepage}
\subsection*{Published neural-model comparison tasks}
\noindent The Neural Latents Benchmark, or NLB, provides shared tasks for comparing models of neural recordings.\par
\end{samepage}
\begin{samepage}
\noindent It informs the choice of comparisons. The paper’s own local scores are not official NLB benchmark scores.\par\nopagebreak[4]

\end{samepage}


\noindent\textit{Current registry lookup labels:} Neural Latents Benchmark; NLB\par

% Stable term ID: published-tests-of-changing-recordings
\begin{samepage}
\subsection*{Published tests of changing recordings}
\noindent FALCON is a published benchmark concerned with neural decoding when recordings change and recalibration is limited.\par
\end{samepage}
\begin{samepage}
\noindent The paper uses related source recordings and its own data splits. Its local muscle-signal scores are not official FALCON scores.\par\nopagebreak[4]

\end{samepage}


\noindent\textit{Current registry lookup labels:} FALCON benchmark; FALCON; FALCON M1; FALCON M1-A\par

% Stable term ID: ranges-that-need-statistical-assumptions
\begin{samepage}
\subsection*{Ranges that need statistical assumptions}
\noindent A confidence interval gives a range for a quantity being estimated. A conformal prediction range tries to contain a future actual result at a stated frequency when its required assumptions hold.\par
\end{samepage}
\begin{samepage}
\noindent The transfer figure shows one fixed run. The paper’s ranges are based on earlier prediction errors, with related errors across time; they do not guarantee how often future true values will fall inside.\par\nopagebreak[4]

\end{samepage}


\noindent\textit{Current registry lookup labels:} Confidence interval and conformal coverage; confidence interval; repeated-seed interval; Prediction range with tested coverage; conformal; conformal coverage; prospective-conformal; distribution-free coverage\par

% Stable term ID: recent-signal-model
\begin{samepage}
\subsection*{Recent-signal model}
\noindent A model component that keeps a learned memory of recent inputs as it predicts the next result.\par
\end{samepage}
\begin{samepage}
\noindent Its formal name is gated recurrent unit, or GRU. This established method uses completed input history; its memory does not establish a universal neural language.\par\nopagebreak[4]

\end{samepage}


\noindent\textit{Current registry lookup labels:} Gated recurrent unit; GRU; compact recurrent model; recurrent network\par

% Stable term ID: record-of-what-was-observed
\begin{samepage}
\subsection*{Record of what was observed}
\noindent A record showing which channels or intervals were actually measured and usable.\par
\end{samepage}
\begin{samepage}
\noindent Missing observations stay missing; they do not become zeros. The second Indy source does not supply separate observation intervals for every recorded unit.\par\nopagebreak[4]

\end{samepage}


\noindent\textit{Current registry lookup labels:} Observation mask; observation interval; invalid support; missing observation\par

% Stable term ID: recorded-neural-events
\begin{samepage}
\subsection*{Recorded neural events}
\noindent A spike is an electrical event identified in a recording. A spike train is the sequence of those events over time.\par
\end{samepage}
\begin{samepage}
\noindent A recorded event count is a whole number tied to one channel and time interval. A model’s average expected count can be a fraction and is a different kind of value.\par\nopagebreak[4]

\end{samepage}


\noindent\textit{Current registry lookup labels:} Spike and spike train; neural event; spike count; spike train\par

% Stable term ID: records-linking-files-checks-and-claims
\begin{samepage}
\subsection*{Records linking files, checks and claims}
\noindent A manifest lists specified files and identities; a receipt records checks; an evidence matrix links a claim to its supporting record.\par
\end{samepage}
\begin{samepage}
\noindent These records also keep failures and withdrawn results visible. A catalog listing does not show that every dataset was acquired or used to train a model.\par\nopagebreak[4]

\end{samepage}


\noindent\textit{Current registry lookup labels:} Manifest, receipt and evidence matrix; manifest; source manifest; audit receipt; processing receipt; evidence matrix; withdrawal record\par

% Stable term ID: relay
\begin{samepage}
\subsection*{Relay}
\noindent Pass focused work between agents, then check their results before accepting it.\par
\end{samepage}
\noindent A manual coordination protocol that passes bounded work and verified evidence between specialist sessions, with one integration owner and an independent acceptance reviewer. Each session follows the currently authorized model, effort and speed routing.\par
\noindent A specialist can work on a bounded question while an independent reviewer checks its evidence before the integration owner accepts it. This is a project coordination workflow, distinct from an established external standard or native Ultra mode. Efficiency and quality claims apply only to tested tasks; universal performance equivalence and a 50 percent weekly saving remain undemonstrated.\par
\begin{samepage}
\noindent\textit{Computational granularity:} finite delegated tasks and evidence artifacts\par\nopagebreak[4]

\end{samepage}


\noindent\textit{Current registry lookup labels:} relay; manual Relay protocol; Relay coordination\par

% Stable term ID: rescaled-muscle-signal
\begin{samepage}
\subsection*{Rescaled muscle signal}
\noindent A processed summary of muscle electrical activity put onto a relative scale.\par
\end{samepage}
\begin{samepage}
\noindent The 16 targets have no physical voltage unit. A stale millivolt label in the source does not make them measured voltage, force or muscle activation.\par\nopagebreak[4]

\end{samepage}


\noindent\textit{Current registry lookup labels:} Normalized EMG envelope; EMG envelope; quantile normalization; dimensionless EMG\par

% Stable term ID: rescaling-using-training-data-only
\begin{samepage}
\subsection*{Rescaling using training data only}
\noindent Learning a rescaling rule from the training portion, then applying the same rule to later values.\par
\end{samepage}
\begin{samepage}
\noindent This avoids using local test values to fit that rule. It cannot undo source processing that already used future samples or the whole recording.\par\nopagebreak[4]

\end{samepage}


\noindent\textit{Current registry lookup labels:} Training-only normalization; scaler; target statistics; normalization\par

% Stable term ID: response-after-an-electrical-stimulus
\begin{samepage}
\subsection*{Response after an electrical stimulus}
\noindent An electrical recording made while observing the response to a stimulus.\par
\end{samepage}
\begin{samepage}
\noindent The human preview is a stimulus-response experiment, not voluntary movement. The stimulus itself can also leave an unwanted signal in the recording; that is an artifact.\par\nopagebreak[4]

\end{samepage}


\noindent\textit{Current registry lookup labels:} Evoked potential and stimulus artifact; evoked stimulation; evoked recording\par

% Stable term ID: results-already-seen-during-development
\begin{samepage}
\subsection*{Results already seen during development}
\noindent Results that have been inspected and can therefore influence later decisions.\par
\end{samepage}
\begin{samepage}
\noindent The human 90\% result and local motor test windows have this limitation. They can be reported, but their later reuse is not a new independent test.\par\nopagebreak[4]

\end{samepage}


\noindent\textit{Current registry lookup labels:} Development evidence; development holdout; pilot; descriptive local test; inspected test\par

% Stable term ID: reusing-learning-in-another-recording
\begin{samepage}
\subsection*{Reusing learning in another recording}
\noindent Trying to use information learned in one setting to improve performance in another.\par
\end{samepage}
\begin{samepage}
\noindent The comparison must state the new-setting labels and the controls. The reported frozen motor transfer and unadjusted human cross-day result do not establish successful biological or cross-species transfer.\par\nopagebreak[4]

\end{samepage}


\noindent\textit{Current registry lookup labels:} Cross-recording transfer; transfer; cross-day transfer; cross-species transfer\par

% Stable term ID: routing-a-muscle-around-a-modeled-obstacle
\begin{samepage}
\subsection*{Routing a muscle around a modeled obstacle}
\noindent A rule bends a modeled muscle route around an obstacle. One obstacle shape, a torus, is shaped like a doughnut.\par
\end{samepage}
\begin{samepage}
\noindent The correction restricts three elbow wrapping rules to a stated attachment segment. It removes reproduced route errors without deleting a muscle or claiming a new anatomical measurement.\par\nopagebreak[4]

\end{samepage}


\noindent\textit{Current registry lookup labels:} Muscle wrapping and torus; wrap; wrapping route; torus; ELBOW\_ULNA; ELBOW\textbackslash{}\_ULNA; off-torus contact\par

% Stable term ID: saved-model-and-matching-output-check
\begin{samepage}
\subsection*{Saved model and matching-output check}
\noindent A checkpoint saves a model’s state. A parity check asks whether another execution produces the specified matching output.\par
\end{samepage}
\begin{samepage}
\noindent Repeating a result needs the same support, ordering, conversions, data splits, settings and body-model version. Passing software checks does not certify biological validity.\par\nopagebreak[4]

\end{samepage}


\noindent\textit{Current registry lookup labels:} Checkpoint and parity; checkpoint; replay/checkpoint parity; model parity; reproduction\par

% Stable term ID: scalp-electrical-recording
\begin{samepage}
\subsection*{Scalp electrical recording}
\noindent Electrodes on the scalp record electrical potential; the method is called EEG.\par
\end{samepage}
\begin{samepage}
\noindent This is an aggregate electrical signal, not a list of individual cortical spikes. The project’s paired preview follows electrical stimulation rather than voluntary motor training.\par\nopagebreak[4]

\end{samepage}


\noindent\textit{Current registry lookup labels:} Electroencephalography; EEG; scalp potentials\par

% Stable term ID: settings-needed-to-repeat-a-run
\begin{samepage}
\subsection*{Settings needed to repeat a run}
\noindent The numbers and choices fixing one training run’s starting point and procedure.\par
\end{samepage}
\begin{samepage}
\noindent One fixed run can be reproduced. It does not show how results vary across repeated starts, animals or future recording days.\par\nopagebreak[4]

\end{samepage}


\noindent\textit{Current registry lookup labels:} Random seed and configuration; seed; single-seed; fixed protocol\par

% Stable term ID: shared-internal-model
\begin{samepage}
\subsection*{Shared internal model}
\noindent A model can learn an internal pattern and reuse that pattern when making several predictions.\par
\end{samepage}
\begin{samepage}
\noindent The paper calls the pattern a latent representation and the reusable model a backbone or trunk. The tested shared model and frozen-transfer candidates were rejected.\par\nopagebreak[4]

\end{samepage}


\noindent\textit{Current registry lookup labels:} Latent representation and backbone; latent model; shared latent; backbone; trunk; shared parameters\par

% Stable term ID: simple-weighted-prediction-model
\begin{samepage}
\subsection*{Simple weighted prediction model}
\noindent A linear model that combines inputs using fitted weights while discouraging excessively large weights.\par
\end{samepage}
\begin{samepage}
\noindent The formal method is ridge regression. It provides a meaningful comparison: a more complex model must add value rather than merely produce a plausible display.\par\nopagebreak[4]

\end{samepage}


\noindent\textit{Current registry lookup labels:} Ridge regression; ridge; linear control; linear model\par

% Stable term ID: simplicity-principle
\begin{samepage}
\subsection*{Simplicity principle}
\noindent Introduce one concept or element at a time, using separate slides or successive reveals according to what makes the relationship easiest to understand.\par
\end{samepage}
\noindent Introduce each new element on its own, then show how the elements connect. Use separate slides or successive reveals to make the logical relationship visible. Simplicity cannot remove scientific qualifications, evidence, source rights or required methods. Actual comprehension and exported readability still need review.\par


\noindent\textit{Current registry lookup labels:} No additional label declared.\par

% Stable term ID: skin-prick-suit
\subsection*{Skin-prick suit}
\noindent A suit shows proposed interface contacts across the body.\par
\noindent The proposed body-side (peripheral) interface, drawn as a suit.\par
\noindent A contextual display label in the user's Interoceptomics visual instructions, not a general clinical synonym. The name establishes no acquired neural data, registration, tissue access, function, safety, clinical use, species or physical dimensions.\par

\noindent\textit{Current registry lookup labels:} skin prick suit; peripheral suit; peripheral prosthesis; peripheral interface\par

% Stable term ID: small-connector-model
\begin{samepage}
\subsection*{Small connector model}
\noindent A smaller fitted mapping that connects one recording to an existing model or interface.\par
\end{samepage}
\begin{samepage}
\noindent The connector can change while the older model stays fixed. That does not establish matching neurons or useful transfer without new-setting calibration.\par\nopagebreak[4]

\end{samepage}


\noindent\textit{Current registry lookup labels:} Adapter; linear adapter; recording adapter\par

% Stable term ID: software-that-calculates-a-computer-body-s-motion
\begin{samepage}
\subsection*{Software that calculates a computer body’s motion}
\noindent MuJoCo calculates how a modeled body moves under forces and control inputs.\par
\end{samepage}
\begin{samepage}
\noindent It is established software reused here. Its availability does not prove that the chosen body model or mapping matches a particular animal’s physiology.\par\nopagebreak[4]

\end{samepage}


\noindent\textit{Current registry lookup labels:} MuJoCo; dynamics engine\par

% Stable term ID: spinal-strip
\subsection*{Spinal strip}
\noindent A strip shows the proposed interface along the modeled spine.\par
\noindent The proposed spinal interface, drawn as a strip along the modeled spine.\par
\noindent A contextual display label in the user's Interoceptomics visual instructions, not a general clinical synonym. The name establishes no acquired neural data, registration, tissue access, function, safety, clinical use, species or physical dimensions. It does not validate a modeled dorsal-surface or tailbud overlay as biological spinal-cord or conus anatomy.\par

\noindent\textit{Current registry lookup labels:} spinal strip; spinal strap; spinal prosthesis; spinal interface\par

% Stable term ID: spine2body
\begin{samepage}
\subsection*{Spine2Body}
\noindent A muscle signal becomes a movement through tendons, joints, contact and load. Measure those steps together, and a failed reach becomes a specific question about activation, coordination or mechanics.\par
\end{samepage}
\noindent The pathway through which spinal output reaches peripheral motor nerves and muscles, producing action through the body's mechanics, while sensory signals return from the body.\par
\noindent Surface EMG measures muscle electrical activity, not peripheral nerve firing or force. Spinal plasticity, momentary reflex modulation and peripheral adaptation remain distinct. Simulated movement is a prediction to compare with independent measurements.\par
\begin{samepage}
\noindent\textit{Organizational level:} whole system; network; signal/feature\par\nopagebreak[4]

\end{samepage}


\noindent\textit{Current registry lookup labels:} spine2body; spine to body; spine-to-body\par

% Stable term ID: stopping-training-under-a-stated-rule
\begin{samepage}
\subsection*{Stopping training under a stated rule}
\noindent Ending training when the chosen development rule says to stop, rather than always using every allowed round.\par
\end{samepage}
\begin{samepage}
\noindent Compared methods have equal round limits but can perform different numbers of actual updates. This is not measured biological learning time.\par\nopagebreak[4]

\end{samepage}


\noindent\textit{Current registry lookup labels:} Early stopping; epoch cap; training update\par

% Stable term ID: structured-task-text-and-supported-instructions
\begin{samepage}
\subsection*{Structured task text and supported instructions}
\noindent JSON stores structured text values. The task compiler converts a supported instruction into a defined program.\par
\end{samepage}
\begin{samepage}
\noindent The compiler chooses available recorded building blocks and limited computer-controller practice. It does not program arbitrary biological skills or supply missing task evidence.\par\nopagebreak[4]

\end{samepage}


\noindent\textit{Current registry lookup labels:} JSON and task compiler; JSON; task JSON; natural-language compiler; supported primitive\par

% Stable term ID: surface-electrical-recording-over-the-spine
\begin{samepage}
\subsection*{Surface electrical recording over the spine}
\noindent Surface electrodes placed over spinal regions record electrical potentials; the method is called ESG.\par
\end{samepage}
\begin{samepage}
\noindent These measurements combine electrical contributions. They are not separately identified spinal-neuron spike trains.\par\nopagebreak[4]

\end{samepage}


\noindent\textit{Current registry lookup labels:} Electrospinography; ESG; surface spinal recording; surface ESG\par

% Stable term ID: system-3-thinking
\begin{samepage}
\subsection*{System 3 thinking}
\noindent System 3 thinking is the partnership; a large movement model is the movement-focused model we aim to build within it.\par
\end{samepage}
\noindent The programme’s human–AI partnership framing in which an external supervisory model helps interpret, train or calibrate task-specific neural/body mappings; the movement-focused model being built is called a large movement model.\par
\noindent Expert footage provides movement examples, but it does not supply the depicted experts’ neural or muscle recordings. An external supervisor would help interpret, train or calibrate task-specific neural/body mappings while the large movement model remains the movement-focused model being built within that partnership. Shaw and Nave use System3 for artificial cognition outside the brain in a different research account. This programme’s task-specific framing establishes no new anatomical brain layer, biological neo-neocortex, third anatomical thinking system or achieved universal nervous-system model.\par
\noindent\textit{Organizational level:} Declared game/model/interface role\par
\noindent\textit{Computational granularity:} Declared task-specific model or interface\par
\begin{samepage}
\noindent\textit{Statistical aggregation:} Validation scope declared per experiment; no physical-scale inference\par\nopagebreak[4]

\end{samepage}
\noindent\textit{Source:} Thinking—Fast, Slow, and Artificial\par
\noindent\textit{Source identifier:} shaw-nave2026\par
\noindent\textit{DOI:} 10.2139/ssrn.6097646\par
\noindent\url{https://papers.ssrn.com/sol3/papers.cfm?abstract_id=6097646}\par
\noindent Prior proposed System3 usage for artificial cognition outside the brain; not an established third anatomical brain system\par
\noindent\textit{Source:} Thinking Fast, Slow, Artificially: AI and Your Brain\par
\noindent\textit{Source identifier:} wharton-authors2026\par
\noindent\href{https://executiveeducation.wharton.upenn.edu/thought-leadership/wharton-at-work/2026/05/thinking-fast-slow-and-artificially/}{\nolinkurl{https://executiveeducation.wharton.upenn.edu/}}\\
\href{https://executiveeducation.wharton.upenn.edu/thought-leadership/wharton-at-work/2026/05/thinking-fast-slow-and-artificially/}{\nolinkurl{thought-leadership/wharton-at-work/2026/05/}}\\
\href{https://executiveeducation.wharton.upenn.edu/thought-leadership/wharton-at-work/2026/05/thinking-fast-slow-and-artificially/}{\nolinkurl{thinking-fast-slow-and-artificially/}}\par
\noindent Primary institutional account of Shaw and Nave System3 research\par


\noindent\textit{Current registry lookup labels:} System three thinking; System3 thinking\par

% Stable term ID: testing-on-an-excluded-day-animal-or-species
\begin{samepage}
\subsection*{Testing on an excluded day, animal or species}
\noindent Checking a model on complete identities that were kept out of the relevant development steps.\par
\end{samepage}
\begin{samepage}
\noindent The paper uses formal names such as leave-session-out or leave-species-out. A good score within one session does not answer these broader tests.\par\nopagebreak[4]

\end{samepage}


\noindent\textit{Current registry lookup labels:} Leave-out generalization; leave-session-out; leave-animal-out; leave-species-out; unseen-organism support\par

% Stable term ID: testing-possible-next-actions-in-the-model
\begin{samepage}
\subsection*{Testing possible next actions in the model}
\noindent Simulating near-future actions before choosing the next command.\par
\end{samepage}
\begin{samepage}
\noindent The copied test states must not change the live model. The native result remains rejected and too slow for real-time use; cached playback is a separate function.\par\nopagebreak[4]

\end{samepage}


\noindent\textit{Current registry lookup labels:} Short-horizon planning and lookahead; horizon controller; lookahead; copied-state rollout; model predictive planning\par

% Stable term ID: the-comparison-needed-to-judge-improvement
\begin{samepage}
\subsection*{The comparison needed to judge improvement}
\noindent A defined comparison that shows whether the proposed method adds useful information or performance.\par
\end{samepage}
\begin{samepage}
\noindent A simple predictor, random fixed model and constant muscle command answer different questions. Producing a plausible animation is not a substitute for beating the relevant comparison.\par\nopagebreak[4]

\end{samepage}


\noindent\textit{Current registry lookup labels:} Baseline and control; simple baseline; tonic baseline; random frozen control; behavior-only predictor\par

% Stable term ID: thought2action
\begin{samepage}
\subsection*{Thought2Action}
\noindent Begin with one intention and one action: reach for the cup. Measure the neural, muscular and mechanical steps between them. The aim is to make that intent-to-action relationship as close to one-to-one as possible within the declared task, person, state, context and tolerated error.\par
\end{samepage}
\noindent A task-conditioned mapping from a participant's stated or inferred intent to a specified action, evaluated through independently measured outcomes.\par
\noindent Intent inferred from EEG is not direct access to private thought. Movement constrains competing explanations; it does not certify an achieved universal bijection or the best possible mapping.\par
\begin{samepage}
\noindent\textit{Organizational level:} whole system; network; signal/feature\par\nopagebreak[4]

\end{samepage}


\noindent\textit{Current registry lookup labels:} thought2action; thought to action; thought-to-action\par

% Stable term ID: thought2world
\begin{samepage}
\subsection*{Thought2World}
\noindent Follow the intention all the way to its consequence, then follow sensation and perception back to the next intention. Measuring the middle makes errors easier to locate and the next experiment easier to define.\par
\end{samepage}
\noindent The composed pathway Thought2Action followed by Action2World: intent is translated into an action, and that action produces an outcome that can be independently checked.\par
\noindent This is a proposed programme pathway, not a claim of unrestricted thought control. It is not silently equivalent to the existing Body2World lookup label, Body2Action or Brain2Motion. External touch and movement tasks alone do not solve internal-body sensing or the full Interoceptomics programme.\par
\begin{samepage}
\noindent\textit{Organizational level:} whole system; network; signal/feature\par\nopagebreak[4]

\end{samepage}


\noindent\textit{Current registry lookup labels:} thought2world; thought to world; thought-to-world\par

% Stable term ID: timing-relationships-between-spikes
\begin{samepage}
\subsection*{Timing relationships between spikes}
\noindent How closely spikes can repeat and how spikes from different recorded features occur together.\par
\end{samepage}
\begin{samepage}
\noindent The project checks predicted average counts. It does not establish these finer timing relationships or a complete account of biological spike noise.\par\nopagebreak[4]

\end{samepage}


\noindent\textit{Current registry lookup labels:} Refractoriness and spike synchrony; refractoriness; joint spike synchrony; biological count-noise model\par

% Stable term ID: touch-brain-area
\begin{samepage}
\subsection*{Touch brain area}
\noindent A brain area that processes body-related sensory information. The source study records from somatosensory cortex, including Area 2.\par
\end{samepage}
\begin{samepage}
\noindent The formal shorthand is S1. The selected model estimates an earlier limb state; it does not recover voluntary motor commands or spinal and nerve recordings.\par\nopagebreak[4]

\end{samepage}


\noindent\textit{Current registry lookup labels:} Somatosensory cortex and Area 2; S1; Area 2; Area2\_Bump; somatosensory\par

% Stable term ID: guide-and-learner
\subsection*{Trainer and trainee}
\noindent The trainer supplies guidance; the trainee’s later control is tested. An experienced trainer helping a novice asks a different question from two novices practicing together.\par
\noindent The proposed comparisons include experienced-to-novice, novice-to-novice, individual practice and matched no-link conditions. Later trainee-alone testing removes trainer commands while retaining any bridge required to express movement.\par

\noindent\textit{Current registry lookup labels:} Tutor and pupil; biological tutor; artificial tutor; live tutoring; tutor withdrawal; experienced-to-novice tutor; novice-to-novice; individual practice; Trainer and learner; biological trainer; artificial trainer; live training; trainer withdrawal; experienced-to-novice trainer; Guide and learner; Guide and trainee\par

% Stable term ID: training-a-fresh-model
\begin{samepage}
\subsection*{Training a fresh model}
\noindent Training a new model for the current recording without importing an older model’s learned main weights.\par
\end{samepage}
\begin{samepage}
\noindent This is the scratch specialist. It performs better than the frozen-transfer candidate in both reported label budgets.\par\nopagebreak[4]

\end{samepage}


\noindent\textit{Current registry lookup labels:} Scratch training; scratch specialist; new-recording scratch specialist\par

% Stable term ID: training-with-neural-samples-without-movement-labels
\begin{samepage}
\subsection*{Training with neural samples without movement labels}
\noindent Using neural samples for a training objective that does not require paired movement labels.\par
\end{samepage}
\begin{samepage}
\noindent These samples still provide information and require computation. The first 5\% label condition permits full training-neural support, so it is not only 28.5 seconds of total observation.\par\nopagebreak[4]

\end{samepage}


\noindent\textit{Current registry lookup labels:} Neural-only training support; neural-only pretraining; unlabeled neural support\par

% Stable term ID: trying-to-move-and-actually-moving
\begin{samepage}
\subsection*{Trying to move and actually moving}
\noindent Trying or being cued to move is different from measuring the movement that occurred.\par
\end{samepage}
\begin{samepage}
\noindent The human labels describe attempted finger movement. They do not independently separate cue from intent or provide executed finger trajectories.\par\nopagebreak[4]

\end{samepage}


\noindent\textit{Current registry lookup labels:} Attempted versus executed movement; attempted finger identity; cued attempted movement; executed kinematics\par

% Stable term ID: typical-error-on-the-original-scale
\begin{samepage}
\subsection*{Typical error on the original scale}
\noindent RMS or RMSE squares the deviations, averages them, and takes a square root so the result has the original unit.\par
\end{samepage}
\begin{samepage}
\noindent An error per joint differs from a combined seven-joint error. An error against a changed command goal also differs from an error against the same recorded reference.\par\nopagebreak[4]

\end{samepage}


\noindent\textit{Current registry lookup labels:} Root mean square error; RMSE; RMS; root-mean-square; tracking error\par

% Stable term ID: unassigned-event-channel
\begin{samepage}
\subsection*{Unassigned event channel}
\noindent A channel collecting events not assigned to the sorted features used by the model.\par
\end{samepage}
\begin{samepage}
\noindent These events are excluded from the 171 sorted features in the second Indy recording. They cannot silently become another identified cell.\par\nopagebreak[4]

\end{samepage}


\noindent\textit{Current registry lookup labels:} Residual or hash cell; residual/hash; hash slot; threshold crossing\par

% Stable term ID: units-for-distance-speed-and-angle
\begin{samepage}
\subsection*{Units for distance, speed and angle}
\noindent Metres, centimetres, millimetres and micrometres measure distance. Metres per second measures speed; degrees and radians measure angle. Radians per second (rad/s) measures the rate of change of angle.\par
\end{samepage}
\begin{samepage}
\noindent One metre is 100 cm or 1000 mm; one mm is 1000 µm. Right-angle and varied-angle maze layouts differ in geometry, not in a validated difficulty multiplier. Unit conversion or a small model error does not prove anatomical or decoded-pose accuracy.\par\nopagebreak[4]

\end{samepage}


\noindent\textit{Current registry lookup labels:} Length, velocity and angle units; m; metre; cm; centimetre; mm; millimetre; µm; micrometre; um; m/s; rad; radian; degrees; °; rad/s; orthogonal; nonorthogonal; varied-angle maze; maze layout\par

% Stable term ID: units-for-pushes-turns-and-electrical-potential
\begin{samepage}
\subsection*{Units for pushes, turns and electrical potential}
\noindent Newtons measure force; newton metres measure turning moment; millivolts measure electrical potential. A dimensionless value has no physical unit.\par
\end{samepage}
\begin{samepage}
\noindent The handle record mixes force and moment units. The rescaled muscle targets are dimensionless despite a stale mV label and cannot be read as force or activation.\par
\end{samepage}


\noindent\textit{Current registry lookup labels:} Force, moment and voltage units; N; newton; N m; N·m; newton metre; mV; millivolt; dimensionless\par

% Stable term ID: units-for-time-and-sample-rate
\begin{samepage}
\subsection*{Units for time and sample rate}
\noindent Seconds measure time; milliseconds are thousandths of a second. Hertz counts repetitions or samples per second.\par
\end{samepage}
\begin{samepage}
\noindent A 250 Hz source sample interval is 4 ms; a 500 Hz simulation interval is 2 ms. These rates differ from playback pace, biological travel time and time spent computing.\par
\end{samepage}


\noindent\textit{Current registry lookup labels:} Time and frequency units; s; second; seconds; ms; millisecond; milliseconds; Hz; hertz; sample rate; 250 Hz; 500 Hz\par

% Stable term ID: unwanted-stimulus-signal-in-a-recording
\begin{samepage}
\subsection*{Unwanted stimulus signal in a recording}
\noindent Part of a recording caused by the stimulus or measuring setup rather than the intended response alone.\par
\end{samepage}
\begin{samepage}
\noindent Stimulus artifacts remain in the human previews. Their presence and the evoked task limit interpretation as voluntary motor evidence.\par\nopagebreak[4]

\end{samepage}


\noindent\textit{Current registry lookup labels:} stimulus artifact; artifact\par

% Stable term ID: use-in-a-new-setting-without-labeled-adjustment
\begin{samepage}
\subsection*{Use in a new setting without labeled adjustment}
\noindent Applying a learned model to a new setting without using its labels to adjust the model.\par
\end{samepage}
\begin{samepage}
\noindent The second M1 model uses explicit new-recording labels. Its test score is a result after calibration, not zero-shot success.\par\nopagebreak[4]

\end{samepage}


\noindent\textit{Current registry lookup labels:} Zero-shot transfer; zero-shot\par

% Stable term ID: speed-estimated-from-earlier-positions
\begin{samepage}
\subsection*{Velocity estimated from earlier positions}
\noindent Two recorded fingertip positions can give an estimate of how fast the fingertip moved and in which direction between the samples. The signed components describe velocity; speed is its non-negative magnitude.\par
\end{samepage}
\noindent A position-derived velocity estimate formed from the difference between current and earlier recorded positions divided by their elapsed time; individual components retain sign. Speed is the magnitude of the velocity vector.\par
\begin{samepage}
\noindent In the separate calibrated motor recording, each interval uses the last available fingertip sample and the actual elapsed time to form a velocity target in m/s. This makes the calculation causal with respect to available position samples, not proof that the tracker supplied them without filtering or delay. Upstream filtering and latency remain unestablished; the derived target is not an original, directly supplied velocity trace.\par\nopagebreak[4]

\end{samepage}


\noindent\textit{Current registry lookup labels:} Backward finite difference; backward difference; finite difference; causally derived velocity; 250 Hz position support; Speed estimated from earlier positions\par

% Stable term ID: speed-of-each-model-movement
\begin{samepage}
\subsection*{Velocity of each model movement}
\noindent The model’s stored velocity for its movement degrees of freedom, including direction.\par
\end{samepage}
\begin{samepage}
\noindent The formal variable name is qvel. Its signed generalized velocities depend on the model’s degrees of freedom; speed is a magnitude. Component units and frames depend on the declared coordinate and model. In MuJoCo, quaternion orientations can make qpos and qvel different lengths, so qvel is not an entry-by-entry derivative of qpos.\par\nopagebreak[4]

\end{samepage}


\noindent\textit{Current registry lookup labels:} qvel; generalized velocity; joint speed; Speed of each model movement\par

% Stable term ID: what-a-model-estimated
\begin{samepage}
\subsection*{What a model estimated}
\noindent A model’s estimate of a defined result, such as speed, from the inputs it was given.\par
\end{samepage}
\begin{samepage}
\noindent An estimate remains an estimate even when it matches the test data well. It does not become a new measurement of the spinal cord, nerves or muscles.\par\nopagebreak[4]

\end{samepage}


\noindent\textit{Current registry lookup labels:} Decoded or inferred prediction; decoded; inferred; reconstructed; prediction\par

% Stable term ID: what-an-instrument-recorded
\begin{samepage}
\subsection*{What an instrument recorded}
\noindent A value produced by a stated measuring method, such as an electrode or a movement tracker.\par
\end{samepage}
\begin{samepage}
\noindent Recorded data can contain noise, filtering or artifacts. Calculating muscle length from tracked markers does not turn that length into a direct tissue measurement.\par\nopagebreak[4]

\end{samepage}


\noindent\textit{Current registry lookup labels:} Recorded measurement; measured; observed; recorded signal\par

% Stable term ID: what-the-computer-body-produced
\begin{samepage}
\subsection*{What the computer body produced}
\noindent A value produced while running a computer model, such as a modeled muscle’s activity.\par
\end{samepage}
\begin{samepage}
\noindent Computer-body movement and forces can be checked for numerical errors. They remain separate from measurements of a living body and from evidence that an organism learned.\par\nopagebreak[4]

\end{samepage}


\noindent\textit{Current registry lookup labels:} Simulated quantity; simulated; numerical model state\par

% Stable term ID: what-was-calculated-from-tracking
\begin{samepage}
\subsection*{What was calculated from tracking}
\noindent A value calculated from measurements together with a body model, such as an estimated joint angle.\par
\end{samepage}
\begin{samepage}
\noindent The study’s joint angles and muscle lengths use tracked movement and a model. They are not direct recordings of muscle commands, tissue length or peripheral nerve spikes.\par\nopagebreak[4]

\end{samepage}


\noindent\textit{Current registry lookup labels:} Model-derived quantity; tracking-derived; source-derived; biomechanical covariate\par

% Stable term ID: where-a-value-came-from
\begin{samepage}
\subsection*{Where a value came from}
\noindent The record of a value’s source and the steps used to produce it.\par
\end{samepage}
\begin{samepage}
\noindent For example, a measured hand position, an angle calculated from that position and a computer-generated muscle state have different origins. The software must preserve those differences.\par\nopagebreak[4]

\end{samepage}


\noindent\textit{Current registry lookup labels:} Provenance; lineage; source integrity\par

% Stable term ID: which-recorded-channel-is-which
\begin{samepage}
\subsection*{Which recorded channel is which}
\noindent The labels and order used to identify neural features within one recording.\par
\end{samepage}
\begin{samepage}
\noindent The labels do not identify the same neurons across days. The project did not accept unverified matches between cells in different recordings.\par\nopagebreak[4]

\end{samepage}


\noindent\textit{Current registry lookup labels:} Source-local unit identity; ordered channel identity; cell correspondence; unit order; neuron index\par

% Stable term ID: which-samples-a-value-depends-on
\begin{samepage}
\subsection*{Which samples a value depends on}
\noindent The input history and exact time interval needed to produce a value or target.\par
\end{samepage}
\begin{samepage}
\noindent A single timestamp does not show all of this support. Overlapping histories and invalid samples near data splits need explicit checks.\par\nopagebreak[4]

\end{samepage}

\begin{samepage}
\subsection*{Composed pathways}
\noindent The proposed intermediary has a positive purpose: understand and measure the processes between intent and action so the task-specific mapping can be made as close to one-to-one as possible. The task, participant, state, context, target and tolerated error must be specified and independently tested. This is an aim for identifiability and control, not an achieved universal bijection or a demonstrated optimum.\par
\end{samepage}
\noindent Form \texorpdfstring{\ensuremath{\rightarrow}}{->} function \texorpdfstring{\ensuremath{\rightarrow}}{->} action and intent \texorpdfstring{\ensuremath{\rightarrow}}{->} action describe connected questions. The following proposed pathways separate the links that would need independent tests before their combined prediction could be trusted.\par
\noindent\textit{Biological pathway:} Brain2Spine \texorpdfstring{\ensuremath{\rightarrow}}{->} Spine2Body; resulting in Brain2Body\par
\noindent Relate brain activity to spinal processing, and spinal/peripheral processing to the body; compose the tested links into an intent-to-body-action pathway.\par
\noindent\textit{Functional pathway:} Thought2Action \texorpdfstring{\ensuremath{\rightarrow}}{->} Action2World; resulting in Thought2World\par
\noindent Relate a declared intention to a measured action, and that action to its measured consequence in the task environment.\par
\noindent\textit{Molecular analogy:} DNA \texorpdfstring{\ensuremath{\rightarrow}}{->} RNA; RNA \texorpdfstring{\ensuremath{\rightarrow}}{->} protein; resulting in DNA \texorpdfstring{\ensuremath{\rightarrow}}{->} protein\par
\noindent The intermediate processes explain how an initial specification reaches its downstream product. This analogy illustrates the need for a measured middle; it does not assert identical biology or a context-free one-to-one relationship.\par
\noindent\textit{System 3 engineering route:} Neural activity \texorpdfstring{\ensuremath{\rightarrow}}{->} Musculoskeletal simulation \texorpdfstring{\ensuremath{\rightarrow}}{->} Programmable model\par
\noindent The proposed engineering route to NeuroTranslation, coordinated by the System3 human–AI supervisory model.\par
\noindent Paired source-bound brain/spinal/peripheral/muscle and sensory/outcome observations constrain mechanics and a task-specific learner. A musculoskeletal simulation represents the biomechanical middle. Model read/write capabilities must be qualified individually.\par
\noindent Intent \texorpdfstring{\ensuremath{\rightarrow}}{->} action \texorpdfstring{\ensuremath{\rightarrow}}{->} outcome returns through sensory feedback to perception and subsequent intent; actual observed/inferred/simulated links remain distinguishable.\par
\noindent Interoceptomics combines interoception and connectomics. Touch, proprioception and external sensorimotor tasks can support this programme; they do not by themselves demonstrate internal-body sensing. Combining the proposed pathways advances a research aim rather than establishing a comprehensive solution.\par
\noindent Brain2Body names the wider programme. Body2World remains a contextual lookup label for Body2Action; Brain2Motion and Body2Action retain their own scopes. These related names do not make their measurement targets equivalent to those of the five composed paths.\par
\noindent The physical spinal pathway remains read-only in this proposal. Each claimed recording or writing capability needs its own independent qualification.\par

\noindent\textit{Current registry lookup labels:} Temporal support and context; support; history window; context; next-target support; overlapping histories\par

% Stable term ID: whole-body-interface
% Canonical status: "user_requested_proposed_programme_role"
\subsection*{Whole Body Interface}
\noindent The hand is one part of the question. Whole Body Interface names the wider ambition to connect declared body systems, while showing which pathways have actually been observed and which remain missing.\par
\noindent A proposed extension of brain–body interfacing toward specified observations and outputs across the brain, spinal, peripheral and body systems.\par
\noindent Whole body states an ambition, not achieved coverage. An interface exchanges specified observations or outputs; the existing model spanning the body is a related modelling ambition, not the same object. State the recorded systems, missing pathways, units, clocks, frames and output routes. Separate datasets can support qualified learning and transfer, but they do not become simultaneous measurements of one body. Physiological input/output direction, software input/output and ordinary sensory feedback remain distinct; no spinal writer or physical actuation is authorized.\par


\noindent\textit{Current registry lookup labels:} No additional label declared.\par

% Stable term ID: xenotransgenicism
% Canonical status: "user_requested_working_label_scope_unsettled"
\subsection*{XenoTransgenesis}
\noindent Which genetic feature would be studied, in which organism, and what would count as the relevant change? XenoTransgenesis is the user’s working label for that proposed question within EndoHumanisation.\par
\noindent A working roadmap name for proposed research involving specified genetic components from another species in a model organism.\par
\begin{samepage}
\noindent XenoTransgenesis is the preferred spelling; XenoTransgenicism remains a historical lookup label. This is not an established replacement name for every form of transgenesis, genome editing or humanisation; use the established term appropriate to any actual manipulation. Neither a genetic change nor its name establishes the intended functional effect or human equivalence. No procedure, approved organism, developmental schedule or experimental permission is defined here.\par\nopagebreak[4]

\end{samepage}


\noindent\textit{Current registry lookup labels:} XenoTransgenicism\par

\subsection*{A cross-species graft}

\noindent A chimera has biological contributions from different sources. A graft is added tissue; when it crosses species, it is a xenotransplant. An organoid is an organized tissue model.\par

\noindent Survival, physical connection, activity and useful function are separate questions. The proposed comparisons do not establish a humanized whole brain or an implemented treatment.\par

\begin{samepage}

\noindent Cross-species-graft component only; not a strict synonym for every chimera, graft or organoid, or achieved ExoHumanisation.\par\nopagebreak[4]

\noindent\textit{Additional lookup label:} XenoTransplantation\par

\end{samepage}

\subsection*{Proposed names and explanatory comparisons}

% Proposal ID: neo-neocortex-inspiration

\subsubsection*{Neo-neocortex}

\noindent A metaphor for an external machine contribution to human cognition; no anatomical addition is claimed.\par

% Proposal ID: alpha-transcriptome-capability-placeholder

\subsubsection*{AlphaTranscriptome}

\noindent AlphaTranscriptome is a proposed name for a desired RNA/cell-state prediction capability. It is not an adopted general-purpose model. AlphaGenome already predicts some expression and splicing~\citep{alphagenomeRetained}; other RNA/cell-state models have bounded domains. These precedents leave whole-cell programmability to be independently demonstrated.\par

% Proposal ID: alpha-cell-supplied-passage-concept

\subsubsection*{AlphaCell}

\noindent A label for integrating cellular layers appears in the supplied passage; no implementation, novelty or product entitlement follows.\par

% Proposal ID: connetype-individual-profile-proposal

\subsubsection*{Connetype}

\noindent A connection profile would describe the selected observations for one individual. Its anatomical coverage, resolution, method and uncertainty must remain visible. The profile does not identify that person’s unique function, thoughts or a fixed biological type.\par

\noindent An individual’s specified neural-connection profile under a declared measurement and representation.\par

\noindent The original three-name line ends in a question mark; all three remain tentative. No formal empirical definition or repeatability threshold is fixed.\par

\noindent The tentative name is Connetype. Connectype remains an unapproved editorial alternative; it is neither a lookup alias nor an adopted spelling.\par

% Proposal ID: interotype-individual-profile-proposal

\subsubsection*{Interotype}

\noindent A heartbeat-related observation can vary with state and measurement conditions. A proposed individual profile must state what was observed and how repeatability would be assessed; it cannot treat touch or proprioception as interchangeable with interoception.\par

\noindent An individual’s profile of specified internal-body sensing observations under declared conditions.\par

\noindent The original three-name line ends in a question mark; all three remain tentative. No formal empirical definition or repeatability threshold is fixed.\par

\noindent Interotype is a tentative name, rather than an established scientific type.\par

% Proposal ID: enactotype-individual-profile-proposal

\subsubsection*{Enactotype}

\noindent A profile of selected actions is bounded by the tasks, observations and conditions used to construct it. Repeatability, state dependence and generalization remain empirical questions; a stable innate action type or inferred unique intent is not established.\par

\noindent An individual’s profile of specified actions and action–outcome observations in declared tasks.\par

\noindent The original three-name line ends in a question mark; all three remain tentative. No formal empirical definition or repeatability threshold is fixed.\par

\noindent Enactotype is a tentative name, rather than an established scientific type.\par

% Proposal ID: form-function-feature-explanatory-lens

\subsubsection*{Form, Function, Feature}

\noindent Start with a structure, ask what it does in the stated setting, then name the observable feature being compared. A hand’s form, its grasping function and a measured property of a particular grasp answer related but different questions.\par

\noindent Feature must identify the actual observed or defined trait in its context; it is not a new universal third biological layer or a substitute for function or action. The existing Form \(\rightarrow\) function \(\rightarrow\) action account and intent–action distinction remain unchanged. Structure constrains function; context, history, selection and co-option prevent a guaranteed one-to-one or teleological chain.\par

% Proposal ID: semiosis-sensorimotor-explanatory-analogy

\subsubsection*{Semiosis as a sensorimotor analogy}

\noindent A drawing can resemble a hand, a pointing finger can indicate an object, and an instruction can follow a learned convention. Such signs can guide action through different relations; they do not assign perception, sensation and action to three exclusive boxes.\par

\noindent Icon, index and symbol describe sign–object relations, while sensation, perception and action name processes. A sign can combine relations. No one-to-one biological correspondence, neural code, private-experience identity or strict isomorphism is established. Outside-in and inside-out perspectives must state whether they concern physiology, modelling or interpretation.\par

% Proposal ID: agent-vs-neurology-stage5-proposal

\subsubsection*{Agent vs Neurology}

\noindent The roadmap ends by asking what an agent could contribute to a defined neurological problem. Agent vs Neurology retains that proposed destination; the specific problem, task and comparison still have to be chosen.\par

\noindent An unresolved Stage 5 research question, not a selected disorder, intervention, diagnostic system or clinical trial.\par

\noindent No specific disease, patient population, endpoint, treatment mechanism, universal cure or successful clinical benefit is inferred. Any future study must define those matters and its ordinary feedback or separately qualified physical interface under the original scientific, participant, rights, safety and independent gates. A prior software, metadata or model result cannot supply this evidence.\par

\noindent Agent vs Neurology names Stage 5 of the proposed roadmap.\par
